% ============================================================================
% Affordable Mistakes -- ECOOP (LIPIcs)
%
% Relationship to the rest of paper/:
%   * paper/tas.tex is the AUTHORITATIVE statement of the BASE discipline's
%     rules.  Section 3 restates notation only; on any divergence tas.tex wins.
%   * paper/skillachievability.tex is the paper on the base discipline.
%   * proof/ mechanizes Sections 4-8 and 11; every result the text marks with
%     \coqok is checked axiom-free by a Print Assumptions harness, and
%     scripts/check_paper_citations.py fails if one is not.
%   * proof/DeviationLayer.v + proof/AUDIT.md is the mechanized audit that
%     refuted the PREDECESSOR draft (modes/taint/grades) and motivated the
%     present reframing; recounted in Section 13 (\S sec:whatchanged).
%
% Before submitting: `make check' must pass (it ties every number below to a
% file in results/ and every \coqok name to a proof).  The `anonymous' class
% option is ON for double-blind; drop it, and fill in \author and \funding,
% for the camera-ready.  The bibliography is main.bib with plainurl (run bibtex).
% ============================================================================
\documentclass[a4paper,UKenglish,cleveref,autoref,thm-restate,anonymous]{lipics-v2021}

\usepackage{mathpartir}
\usepackage{stmaryrd}
\usepackage{booktabs}
\usepackage{pifont}
\usepackage{tikz}
\usetikzlibrary{arrows.meta,positioning,calc}

\definecolor{steel}{RGB}{46,74,107}
\definecolor{boxbg}{RGB}{243,244,246}
\definecolor{okgreen}{RGB}{0,158,115}
\definecolor{verm}{RGB}{213,94,0}

% -------------------------------------------------------- base macros (subset)
% Same meanings as in tas.tex / skillachievability.tex; only the ones this
% paper uses are kept.
\newcommand{\CapS}{\mathsf{Cap}}
\newcommand{\Pred}{\mathsf{Pred}}
\newcommand{\Var}{\mathsf{Var}}
\newcommand{\Lab}{\mathsf{Lab}}
\newcommand{\pre}{\mathit{pre}}
\newcommand{\eff}{\mathit{eff}}
\newcommand{\Add}{\mathit{Add}}
\newcommand{\Del}{\mathit{Del}}
\newcommand{\apply}{\mathrm{apply}}
\newcommand{\id}[1]{\texttt{#1}}
\newcommand{\coind}{\mathrel{::=^{\mathit{coind}}}}
\newcommand{\MM}{\mathbb{M}}
\newcommand{\ppar}{\;\|\;}
\newcommand{\End}{\mathsf{End}}

\newtheorem{assumption}[theorem]{Assumption}

% LIPIcs already provides: theorem, lemma, corollary, proposition, claim,
% definition, example, remark (bold-triangle style), and proof.

\emergencystretch=3em
% Line numbers stay on (the LIPIcs default): reviewers cite them.

% ---------------------------------------------------------- extra macros (v2)
\newcommand{\Haz}{\mathcal H}
\newcommand{\safeT}[1]{\vdash_{#1}}
\newcommand{\safeRr}[1]{\vdash^{\nu}_{#1}}
\newcommand{\Rob}{\textsf{Robust}}
\newcommand{\reachh}[1]{\rightsquigarrow^{\Haz}_{#1}}
\newcommand{\reachg}[1]{\rightsquigarrow^{\varphi}_{#1}}
\newcommand{\Ben}{\textsf{Benign}}
\newcommand{\Fut}{\textsf{Futile}}
\newcommand{\Cat}{\textsf{Catastrophic}}
\newcommand{\coqok}[1]{\textcolor{okgreen}{\ding{51}}\,\texttt{\small #1}}

\newcommand{\ktol}{$k$-misselection tolerance}
\newcommand{\Ktol}{$k$-misselection-tolerant}
\newcommand{\kstar}{k^{*}}

\title{Affordable Mistakes: Severity-Aware Multiparty Session Types for
Participants that Choose Wrongly}
\titlerunning{Affordable Mistakes}
\author{Anonymous Author}{Anonymous Affiliation, Anonymous City, Anonymous Country}{anon@anon.org}{}{}
\author{Anonymous Author}{Anonymous Affiliation, Anonymous City, Anonymous Country}{anon@anon.org}{}{}
\authorrunning{Anonymous author(s)}
\Copyright{Anonymous author(s)}
\ccsdesc[500]{Theory of computation~Type structures}
\ccsdesc[500]{Theory of computation~Process calculi}
\ccsdesc[300]{Software and its engineering~Software verification}
\keywords{Session types, Multiparty, Behavioural types, Safety, Irreversibility,
Fault tolerance, Mechanized metatheory}
\category{}\relatedversion{}
\supplement{Coq development (265 theorems and 2 constructors, axiom-free, Coq 8.18), the
\texttt{skillc severity} implementation, the severity benchmark, the ported real
protocols, and the evaluation scripts and recorded runs; anonymized for submission.}
\EventEditors{}\EventNoEds{0}
\EventLongTitle{European Conference on Object-Oriented Programming (ECOOP 2027)}
\EventShortTitle{ECOOP 2027}
\EventAcronym{ECOOP}\EventYear{2027}\EventDate{June 28--July 2, 2027}
\EventLocation{Turin, Italy}\EventLogo{}
\SeriesVolume{0}\ArticleNo{0}

\begin{document}
\maketitle

\begin{abstract}
Multiparty session types guarantee that well-typed participants interact correctly, on the
hypothesis that each participant \emph{is} its type. A participant that was instructed rather
than compiled --- a language-model agent --- breaks that hypothesis in one specific way: at a
branch where it should select $A$ it may select $B$. Forbidding this is impossible, and would
be useless if it were possible. We change what the discipline is for: it does not prevent wrong
choices, it prevents catastrophic ones and reports the rest as risk.

The fault model is the participant's own selection at a branch, defined through the guards that
asserted global types already carry. Cascading is a budget $k$ of wrong selections, with no
probabilities anywhere. Budgeted deviation is itself old --- reactive synthesis and
robustification have owned it for a decade. What is new is carrying a quantity a participant
\emph{spends} across a reduction relation and into typed programs: a session typed against a
$k$-tolerant protocol by the base discipline's conformance judgment is hazard-free on every run
of cost at most $k$, budgets distribute over participants, and the guarantee survives recursion
and every permutation of independent nodes. One syntax-directed rule characterizes the condition
exactly; each wrong selection is classified \Ben, \Fut{} or \Cat{} by outcome against a
STRIPS-style world, and these are three consecutive intervals of budgets, not an arbitrary
carve-up. The condition composes against an interface rather than a protocol's internals, and is
decidable for regular protocols. Everything is mechanized in Coq, axiom-free; the tool's own
search is not extracted, but on the boolean fragment its verdict is cross-checked against a
kernel that is.

The evaluation asks whether any of this matters to a real agent. On eighteen protocols
ported from public agent configurations and customer-service policies, each pre-registered
before the tool ran, every verdict matched the expectation and every protocol is $0$-tolerant,
for one reason: its points of no return sit behind rules the policy states and no tool
enforces. Moving those rules into tool preconditions, where a tool can check them, makes $9$ of
the $18$ tolerant at every tested budget; in the other $9$ the rule left is a user confirmation
that no tool can check, and the analysis' narrowing repair makes all $18$ tolerant. In $3500$ runs of two models of one
vendor, agents shown true facts entered a \Cat{} branch of the real protocols in $68$ of $720$
runs, and in none behind a gate built from the classification; agents told by a requester
that every precondition held completed a catastrophe in $53$ of $180$ runs, and in none behind
the gate. The limits are plain: on our own seventeen protocols, with true facts, no agent erred
at all, and the false note is a stress condition of our making.
\end{abstract}


\section{Introduction}\label{sec:intro}

Multiparty session types (MPST) guarantee communication safety, session
fidelity and deadlock-freedom for systems whose participants \emph{are} their
local types~\cite{honda-mpst}: the code was checked, so every run is a
compliant run. A participant that was instructed rather than compiled breaks
that hypothesis in a specific way. At a choice point it may take branch $B$
where the situation demanded branch $A$ --- not by violating the protocol
(both branches are in the type) but by judging the situation wrongly. This is
the ordinary condition of a language-model agent occupying a role, and it is
why such agents are employed: their latitude.

Nothing here is specific to language models. The hypothesis fails for every
participant instructed rather than compiled: an operator following a runbook, a
third-party service whose behaviour is a contract rather than a checked
artifact, a learned controller. Language-model agents make the question urgent
and let our evaluation drive it at scale, but the discipline asks only that the
participant choose, not how.

The tempting response is to make the type system prevent it. That is doubly
wrong: impossible, because the judgement is the participant's, not the
protocol's; and undesirable, because a participant that may never choose
wrongly is one that may never choose. What is needed is \emph{triage}.

\subparagraph*{The reframing.}
Do not ask whether the participant complied. Ask \emph{what a wrong choice
costs}. Three answers matter, and the middle one is what existing disciplines
cannot express (Figure~\ref{fig:idea}):
\begin{itemize}
\item \Ben{} --- a detour; the goal is still reachable.
\item \Fut{} --- the goal is lost, but nothing is harmed. \emph{Failure is not
  disaster.}
\item \Cat{} --- a hazard becomes reachable.
\end{itemize}
A system that blocks \Fut{} is useless; one that permits \Cat{} is dangerous.
Separating them is the whole job. Thesis: \emph{session types say what may
happen; this says what you can survive.}

\subparagraph*{What is and is not new.}
The ingredients are old; the discipline is not. Branch guards on global types
are design-by-contract MPST~\cite{bocchi-concur10,gheri-ecoop22,vassor-ecoop24},
and guard violation on selection is the alarm condition of the monitoring
line~\cite{bocchi-forte13}. Whether a goal remains reachable is dead-end
detection in planning~\cite{steinmetz-aaai16,hoffmann-ecai14}; whether an
effect can be undone is undoability checking~\cite{daum-icaps16,morak-kr25}.
Bounded fault tolerance is classical in
planning~\cite{aineto-ecai23,domshlak-icaps13} and in MPST with crash-stop
failures~\cite{barwell-concur22,barwell-ecoop23}. A syntax-directed system
equivalent to a model checker is a known template~\cite{naik-palsberg,kobayashi-ong}.
We claim none of these. What we claim is C1--C5 below, and nothing wider.

We are careful with the word \emph{type system}. The condition is on global
types, as well-assertedness~\cite{bocchi-concur10} is; participants are typed by
the base discipline's conformance judgment, unchanged. The type-theoretic
content is the theorem coupling the two (Section~\ref{sec:bridge}).

\subparagraph*{Contributions.}
Throughout, \coqok{name} marks a result proved in the accompanying Coq
development and names its identifier there.
\begin{itemize}
\item[\textbf{C1}] \emph{A budget a typed participant spends, carried to
  programs} (\S\ref{sec:mis}, \S\ref{sec:bridge}). A session typed against a
  $k$-tolerant protocol is hazard-free on every run whose misselection cost is
  at most $k$ (\coqok{bridge\_run}); the budget distributes over participants,
  so a global allowance is a contract that can be split and checked per role
  (\coqok{budget\_distributes}); the guarantee survives recursion
  (\coqok{bridge\_mu}) and every permutation of independent nodes, with the
  side conditions discharged syntactically (\coqok{bridge\_interleaved},
  \coqok{strips\_swap\_act}). Budgeted deviation is not ours
  (\S\ref{sec:related}); \emph{carrying a quantity a participant spends
  across a reduction relation, into typed programs}, is what no prior
  discipline does.
\item[\textbf{C2}] \emph{The condition that makes C1 checkable}
  (\S\ref{sec:typing}): one syntax-directed rule, sound and complete for
  $k$-tolerance (\coqok{TC\_exact}), which is exactly non-reachability in a
  product graph (\coqok{TC\_regular}). Read coinductively it is the same
  condition on recursive protocols (\coqok{TR\_exact}, \coqok{TC\_is\_TR}), so
  the rule, the decision procedure and the bridge are about one object at both
  layers; and the tolerance degree $\kstar$, the largest budget a protocol
  survives, is principal: the budgets it tolerates are exactly those
  $\le\kstar$ (\coqok{principal\_characterises}).
\item[\textbf{C3}] \emph{An ordered diagnosis} (\S\ref{sec:severity}): every
  wrong choice is \Ben, \Fut{} or \Cat, and these are not an arbitrary carve-up
  --- the class of a residual is monotone in the budget along the chain
  $\Fut<\Ben<\Cat$ (\coqok{severity\_monotone\_in\_budget}), so the three
  classes are consecutive intervals of budgets and $\kstar$ is a genuine
  threshold (\coqok{tolerance\_degree\_is\_a\_threshold}). Budgeted
  reachability is exactly ``some run with at most $k$ wrong choices reaches a
  hazard'' (\coqok{reach\_mu\_iff\_run}), so the classes are statements about
  executions. \Ben{} is a possibility, not a guarantee, and the guarantee is
  separated from it and shown strictly stronger (\coqok{robust\_benign},
  \coqok{benign\_is\_not\_robust}).
\item[\textbf{C4}] \emph{Decidability and modularity} (\S\ref{sec:regular},
  \S\ref{sec:why}): the $\mu$-unfolding closure is finite
  (\coqok{cands\_closed}) and the decision procedure is proved sound and
  complete (\coqok{decide\_mu\_correct}); sequential composition is typed
  against an interface that carries the remaining budget, not the internals
  (\coqok{TC\_seq\_interface}), and that interface may be projected onto the
  cone of influence of what remains (\coqok{interface\_projection}), which is
  what makes modular checking cheap rather than merely decomposable. Four
  repairs are proved sound, two exactly, and the condition is a congruence, so
  they apply where the tool applies them (\coqok{safeT\_congruence}).
\item[\textbf{C5}] \emph{Evidence} (\S\ref{sec:eval}): an implementation whose
  verdict agrees with a kernel extracted from the proof on 500 of 500 random
  protocols; a severity benchmark of our own; eighteen pre-registered protocols
  ported from public agent configurations and customer-service policies, whose
  hazards are exactly the rules no tool enforces, and a repair that enforces them
  and makes 9 of the 18 tolerant at every tested budget; and 3500 live-agent runs
  in which a gate built from the classification removed every entry into a \Cat{}
  branch and every catastrophe, its null result and stress condition stated.
\end{itemize}
Every theorem above is checked for vacuity, because the predecessor of this
discipline was withdrawn when its own audit found it had none
(\S\ref{sec:whatchanged}): conforming sessions are exhibited for the protocols
the paper reasons about and for their repairs, the guard repair's runtime is
built and shown to have the worlds its theorem needs, and the cone-of-influence
and congruence theorems are instantiated where their hypotheses could have been
satisfiable nowhere (\coqok{cone\_is\_not\_degenerate},
\coqok{congruence\_is\_not\_degenerate}), and \Rob{}, the class the paper
says is strictly stronger than \Ben{}, is shown non-empty
(\coqok{Ggood\_is\_robust}).

A design expectation was refuted by the mechanization and a benchmark
finding refines it: tolerance is anti-monotone in the capability context for
a fixed hazard (\coqok{tolerance\_antitone\_in\_ctx}), yet a new tool that
re-establishes a lost atom can \emph{remove} a point of no return
(\S\ref{sec:eval}).

\begin{figure}[t]
\centering
\resizebox{\textwidth}{!}{\begin{tikzpicture}[>=Stealth, font=\small,
  nd/.style={circle,draw,fill=white,inner sep=1.5pt,minimum size=4.5mm},
  lbl/.style={font=\scriptsize},
  goal/.style={draw=okgreen,thick,fill=okgreen!10,rounded corners,inner sep=3pt},
  dead/.style={draw=gray,thick,fill=gray!12,rounded corners,inner sep=3pt},
  haz/.style={draw=verm,thick,fill=verm!10,rounded corners,inner sep=3pt}]
% world
\node[draw,rounded corners,fill=blue!6,inner sep=3pt,align=center,font=\scriptsize] (W) at (0,1.55) {world $W$\\$\id{fare}=620$};
\node[nd] (C) at (0,0) {};
\node[lbl,left=1pt of C] {$\mathsf p{\to}\mathsf q$};
\draw[->] (W) -- (C);
% intended branch
\node[nd] (A) at (3.2,1.25) {};
\draw[->,thick,okgreen] (C) -- node[above,sloped,lbl,pos=.5] {$\id{stop}[\id{fare}{>}500]$} (A);
\node[lbl,okgreen,above=1pt of A] {\textbf{intended}: guard holds};
\node[goal] (GA) at (7.4,1.25) {goal $\varphi$};
\draw[->,okgreen] (A) -- node[above,lbl] {compliant run} (GA);
% misselection
\node[nd] (B) at (3.2,-1.25) {};
\draw[->,thick,verm,dashed] (C) -- node[below,sloped,lbl,pos=.5] {$\id{go}[\id{fare}{\le}500]$} (B);
\node[lbl,verm,below=1pt of B,xshift=-6mm] {\textbf{misselection}: guard fails};
% three fates
\node[goal] (F1) at (7.4,-0.35) {goal $\varphi$};
\draw[->,okgreen] (B) to[bend left=22] node[above,lbl,pos=.6] {detour} (F1);
\node[dead] (F2) at (7.4,-1.55) {goal lost, nothing harmed};
\draw[->,gray] (B) -- (F2);
\node[haz] (F3) at (7.4,-2.85) {hazard $\Haz$};
\draw[->,verm] (B) -- node[below,sloped,lbl,pos=.38] {$\id{purchase}@\mathsf q$} (F3);
% point of no return marker on the catastrophic edge
\draw[verm,very thick] (5.55,-2.45) -- (5.85,-1.95);
\node[lbl,verm,align=center] at (5.9,-3.35) {point of\\no return};
% severity labels
\node[lbl,okgreen,right=3pt of F1] {\Ben};
\node[lbl,gray,right=3pt of F2] {\Fut};
\node[lbl,verm,right=3pt of F3] {\Cat};
\node[lbl,gray] at (11.35,-1.55) {$\neg\Diamond\varphi\ \wedge\ \neg\Diamond\Haz$};
\node[lbl,okgreen] at (11.35,-0.35) {$\Diamond\varphi\ \wedge\ \neg\Diamond\Haz$};
\node[lbl,verm] at (11.35,-2.85) {$\Diamond\Haz$};
\end{tikzpicture}
}
\caption{The idea, on the running example of \S\ref{sec:overview}. At a guarded
choice in world $W$ the branch whose guard holds is intended; taking another is
a misselection. Its severity is a property of the residual, and the three fates
drawn are the only three: the goal still reachable (\Ben), the goal lost with
no hazard reachable (\Fut), or a hazard reachable (\Cat) --- here by crossing a
point of no return, an irreversible capability after which the goal cannot be
recovered. In $G_{\mathsf{bad}}$ the misselected branch takes the third.}
\label{fig:idea}
\end{figure}

\section{Overview}\label{sec:overview}

\subparagraph*{Running example.}
A planner $\mathsf p$ and a worker $\mathsf q$ holding a purchase tool. Atoms
$\id{verified}$, $\id{booked}$, $\id{funds}$; $\id{verify}$ adds the first,
$\id{purchase}$ adds the second and deletes the third, and no tool restores
$\id{funds}$. The protocol offers two branches:
\[
G_{\mathsf{bad}} \;=\;
\mathsf p{\to}\mathsf q : \bigl\{\;
  \id{safe}.\ \id{verify}@\mathsf q.\ \id{purchase}@\mathsf q.\ \End
  \;,\;\;
  \id{fast}.\ \id{purchase}@\mathsf q.\ \End \;\bigr\}
\]
with goal $\varphi = \id{booked}\wedge\id{verified}$.
Under the default instantiation (\S\ref{sec:mis}) $\id{safe}$ is intended
because the goal survives it and $\id{fast}$ is a misselection because it does
not. The tool reports, and the Coq development proves for the same instance:
$G_{\mathsf{bad}}$ is $0$-tolerant (\coqok{Gbad\_is\_0\_tolerant}) --- if the
participant never misselects, nothing bad happens, which is where classical
MPST stops; it is \emph{not} $1$-tolerant (\coqok{Gbad\_not\_1\_tolerant}) ---
one wrong choice purchases unverified, and $\id{purchase}$ is the point of no
return; and narrowing away $\id{fast}$ is tolerant at every $k$
(\coqok{Ggood\_is\_k\_tolerant}). The headline number is the \emph{tolerance
degree} $\kstar$, the largest $k$ the protocol survives.

For a protocol with two decision points the tool's output is what an engineer
needs (Section~\ref{sec:eval}, \texttt{migration\_backup}):
{\footnotesize
\begin{verbatim}
migration_backup: tolerance degree k* = 1
  choice nodes 4, branches 6, irreversible ['drop_old']
  /choice@ops#0
    backup             Benign        intended
    skip_backup        Catastrophic  misselection  PNR=drop_old
  /choice@ops#0/backup/choice@ops#2
    drop_old           Benign        intended
    keep_old           Benign        intended
  /choice@ops#0/skip_backup/choice@ops#1
    drop_old           Catastrophic  misselection  PNR=drop_old
    keep_old           Benign        intended
  first catastrophe within budget 2:
    choose/skip_backup > act/migrate > choose/drop_old > act/drop_old
  repair (narrow):
    remove skip_backup at /choice@ops#0
    remove drop_old at /choice@ops#0/skip_backup/choice@ops#1
  bystander interleavings: exact (0 independent pairs)
\end{verbatim}}
``Tolerant of one wrong choice; two wrong choices reach data loss, here is the
path and the action that burns the bridge'' is actionable. ``Impossible'' is
not.

\begin{figure}[t]
\centering
\resizebox{\textwidth}{!}{\begin{tikzpicture}[>=Stealth,font=\small,
  box/.style={draw,rounded corners,fill=white,align=center,inner sep=3pt,minimum height=7mm,font=\scriptsize},
  lane/.style={font=\scriptsize\sffamily\bfseries,text=steel},
  outbox/.style={draw=okgreen,fill=okgreen!10,rounded corners,align=center,inner sep=2.5pt,font=\scriptsize}]
% lanes
\fill[boxbg] (-0.2,-3.5) rectangle (3.4,1.95);
\fill[blue!5] (3.9,-3.5) rectangle (10.6,1.95);
\fill[boxbg] (11.0,-3.5) rectangle (14.2,1.95);
\node[lane] at (1.6,1.7) {AUTHOR TIME \ (untrusted)};
\node[lane] at (7.25,1.7) {COMPILE TIME \ (trusted, deterministic)};
\node[lane] at (12.6,1.7) {RUN TIME};
% lane 1
\node[box] (skill) at (1.6,0.85) {\texttt{SKILL.md}\\natural-language skill};
\node[box,dashed] (comp) at (1.6,-0.65) {compaction\\LLM or deterministic};
\node[box] (pack) at (1.6,-2.3) {pack $\langle\Gamma,\ G[\psi],\ \varphi,\ \Haz\rangle$\\guarded choice, goal, hazard};
\draw[->] (skill) -- (comp); \draw[->] (comp) -- (pack);
% trust boundary
\draw[verm,thick,dashed] (3.65,-3.5) -- (3.65,1.95);
\node[verm,font=\scriptsize,rotate=90] at (3.5,0.3) {trust boundary};
% lane 2
\node[box] (gate) at (5.0,0.55) {schema\\gate};
\node[box,minimum width=44mm] (cls) at (8.05,0.55) {\textbf{severity classifier}\\$\forall$ node $\times$ abstract world $\times$ budget $b$:\\$(G_i,W)\reachg{b}$?\quad $(G_i,W)\reachh{b}$?};
\node[draw=okgreen,fill=white,rounded corners,font=\scriptsize,inner sep=2pt] at (6.9,1.45) {\ding{51}\,sound + complete (Coq)};
\node[outbox] (o1) at (5.3,-1.45) {severity\\per branch};
\node[outbox] (o2) at (7.25,-1.45) {tolerance\\degree $k^{*}$};
\node[outbox] (o3) at (9.4,-1.45) {witness +\\PNR action};
\node[outbox,minimum width=40mm] (o4) at (7.6,-2.75) {repairs: guard \,/\, narrow \,/\, reorder};
\draw[->] (pack.east) -- ++(0.35,0) |- (gate.west);
\draw[->] (gate) -- (cls);
\draw[->] (cls.south) ++(-2.2,0) -- (o1.north);
\draw[->] (cls.south) ++(-0.5,0) -- (o2.north);
\draw[->] (cls.south) ++(1.6,0) -- (o3.north);
\draw[->] (o1.south) ++(0.6,0) -- (o4.north -| 6.4,0);
\draw[->] (o3.south) ++(-0.6,0) -- (o4.north -| 9.2,0);
% lane 3
\node[box] (rt) at (12.6,0.55) {gated runtime\\monitor / gate / scheduler};
\node[box] (ag) at (12.6,-1.75) {LLM agents $\mathsf p,\mathsf q,\dots$\\may misselect a branch};
\draw[->] (cls.east) -- node[above=5pt,font=\scriptsize,align=center] {certified\\pack, $k^{*}$} (rt.west);
\draw[<->] (rt) -- node[right,font=\scriptsize,align=left] {attempts\\gated} (ag);
\end{tikzpicture}
}
\caption{Architecture. Compaction from prose to a pack is untrusted and may
use a language model; nothing beyond the trust boundary does. The severity
classifier runs the discipline's existing reachability search pointwise at
every branch of every choice, at every abstract world and budget, and emits a
per-branch severity, the tolerance degree $\kstar$, a witness naming the
point-of-no-return action, and the narrowing repair. The certified pack drives
the gated runtime that refuses out-of-protocol attempts.}
\label{fig:arch}
\end{figure}

\subparagraph*{Architecture.}
Figure~\ref{fig:arch}: a compile-time check on the pack, before any
participant runs, with no language model in the trusted core --- the base
discipline's trust boundary. The runtime gate is inherited unchanged: it
refuses \emph{out-of-set} attempts (labels the protocol does not offer) before
delivery, so those change no state. The deviation this paper analyses is
\emph{in-set}: a branch the type permits and the situation forbids.

\section{The Base Discipline}\label{sec:base}

We build on a goal-marked, capability-guarded synchronous multiparty
discipline, stated here in the detail this paper uses. Predicates
$p\in\Pred$, numeric variables $x\in\Var$, roles $\mathsf p,\mathsf q$,
capabilities $a\in\CapS$, labels $\ell\in\Lab$. The state logic is
quantifier-free linear integer arithmetic with uninterpreted booleans; a world
$W=\langle B,N\rangle$ valuates it, and the checker's decidable fragment
abstracts predicates concretely and numerics symbolically, widening at loop
heads so the reachable world set is finite. A capability context maps each
possessed tool to a guarded STRIPS effect
$\Gamma(a)=\langle\pre_a,\eff_a\rangle$, where $\eff_a$ adds and deletes
predicates and assigns numerics. Firing is \textsc{World-Act}:
$\langle W\rangle\,a\,\langle W'\rangle$ iff $W\vDash\pre_a$ and
$W'\in\apply(\eff_a,W)$. We write $\varphi$ for the goal and $\Haz$ for the
hazard, both predicates over worlds. Global types are coinductive regular trees
\[
G \;\coind\; \mathsf{p}{\to}\mathsf{q}:\{\ell_i.\,G_i\}_{i\in I}
  \ \mid\ a@\mathsf{p}.\,G \ \mid\ \checkmark\varphi.\,G \ \mid\ \End
\]
and sessions $\MM=\mathsf p_1[P_1]\ppar\cdots$ of processes
$P\coind\mathbf 0\mid\mathsf p!\{\ell_i.P_i\}\mid\mathsf p?\{\ell_i.P_i\}\mid a.P$
are typed \emph{directly} against $G$ by the coinductive judgment
$G\vdash(\MM;W)$ --- no local types, projection, or separate
subtyping~\cite{jongmans-ecoop23,castro-popl26} --- whose \textsc{T-Comm}
requires the sender to offer exactly the protocol's labels, the receiver at
least them, and every branch to check against the same residual. A \emph{pack}
$\langle\Gamma,G,\varphi,\Haz\rangle$ is the checked object: a capability
context, a global type, a goal and a hazard. Achievability is existential
may-reachability of a goal-satisfying configuration; the base metatheory
(subject reduction, session fidelity, refutation soundness) is mechanized for
the finite fragment. This paper adds one annotation to choice, one counter to
the semantics, and one well-formedness condition.

\section{Misselection and the Budget}\label{sec:mis}

\begin{definition}[Guarded choice]\label{def:guarded}
A choice node carries, per branch, the predicate that makes that branch the
intended one: $\mathsf p{\to}\mathsf q:\{\ell_i[\psi_i].\,G_i\}_{i\in I}$. In
world $W$, branch $i$ is \emph{intended} if $W\vDash\psi_i$; taking a branch
with $W\nvDash\psi_i$ is a \emph{misselection}.
\end{definition}

The syntax is that of asserted global types~\cite{bocchi-concur10}, whose
well-formedness conditions guarantee that a compliant participant can always
pick a satisfied branch. We keep the syntax and drop the guarantee: the guards
are not a constraint on an implementation but the \emph{definition of wrong}.

\subparagraph*{Default instantiation.}
Guards need not be written. When a branch carries none, the implementation
uses a default derived from the goal and the hazard (\S\ref{sec:severity}).
At a choice in world $W$:
(i)~if the goal is reachable from the residual of some branch, a branch is
intended iff the goal is reachable from its residual;
(ii)~otherwise --- the goal is already lost --- a branch is intended iff the
hazard is \emph{avoidable} from its residual: some maximal run from it (one
that ends, blocks, or returns to a configuration it has visited) fires no
hazard;
(iii)~if no branch avoids the hazard either, every branch is intended.
Before the point of no return the right choice keeps the goal; after it, the
right choice limits harm; a choice with no right answer is free. The budget
therefore counts wrong choices, not the depth of the protocol after the task
was lost. Avoidability is decided on the finite unfolding graph the decision
procedure already explores (\S\ref{sec:regular}), following every branch of
later choices. The recursive alternative, which would follow only the
\emph{intended} branches of later choices, makes the guards of a loop body
depend on themselves; we use the simpler rule, which differs from it only when
an explicit guard further down forbids the one hazard-free continuation (or,
with environment choices, which avoidability takes adversarially, when an
answer of the environment makes the goal reachable again).
Explicit guards override the default where the intended branch is a matter of
policy rather than reachability (Section~\ref{sec:eval} uses both).

\begin{definition}[Budgeted reachability]\label{def:budget}
$(G,W)\reachh{b}$ holds when a hazard state is reachable from $(G,W)$ using at
most $b$ misselections (Coq: \texttt{reach\_haz}): the head rules of the base
semantics, plus
\begin{mathpar}
\inferrule*[left=RH-Comm-Ok]
  {\ell_i[\psi_i].G_i \in \mathit{brs} \\ W\vDash\psi_i \\ (G_i,W)\reachh{b}}
  {(\mathsf p{\to}\mathsf q:\mathit{brs},W)\reachh{b}}
\and
\inferrule*[left=RH-Comm-Dev]
  {\ell_i[\psi_i].G_i \in \mathit{brs} \\ W\nvDash\psi_i \\ (G_i,W)\reachh{b}}
  {(\mathsf p{\to}\mathsf q:\mathit{brs},W)\reachh{b{+}1}}
\end{mathpar}
Compliant progress is free; a misselection costs one unit. A choice
$\mathsf e{:}\{\ell_i.G_i\}$ may instead be controlled by the environment ---
what the user answers, what the records show; one authored and twelve real protocols use it
(\S\ref{sec:eval}). It is no agent decision, so by \textsc{RH-Env} the adversary
takes any branch at no cost: $(G_i,W)\reachh{b}$ gives
$(\mathsf e{:}\{\ell_i.G_i\},W)\reachh{b}$ (\coqok{env\_reach\_iff}).
\end{definition}

\begin{definition}[\ktol{}]\label{def:ktol}
$G$ is \emph{\Ktol{}} for hazard $\Haz$ under $\Gamma$ from $W_0$ if
$\neg\,(G,W_0)\reachh{k}$.\footnote{We avoid two established names. In MPST,
\emph{$k$-safety} is a buffer bound~\cite{lange-yoshida-cav19}; in
hyperproperties it is a property of $k$ traces~\cite{clarkson-schneider}.
\emph{$k$-resilient} plans survive $k$ action failures~\cite{aineto-ecai23};
our fault is the participant's own selection and our target is a hazard, not
the goal.} $0$-tolerance is classical MPST safety.
\end{definition}

The budget is possibilistic: a design parameter of the deployment, not a
measurement of the participant. Every theorem below holds for every
participant that misselects at most $k$ times, however it does so.

\section{Severity and the Point of No Return}\label{sec:severity}

Let $\varphi$ be the goal and $\Haz$ the hazard; $\reachg{b}$ is budgeted goal
reachability.

\begin{definition}[Severity]\label{def:sev}
For a residual $(G,W)$ at budget $b$: \Cat{} if $(G,W)\reachh{b}$; otherwise
\Fut{} if $\neg(G,W)\reachg{b}$; otherwise \Ben{}.
\end{definition}

The trichotomy is not an arbitrary carve-up of two predicates. Both
reachability questions are monotone in the budget
(\coqok{reach\_mono\_budget}), and that makes the classes \emph{ordered}.

\begin{assumption}[Decidable questions]\label{ass:dec}
Where a result needs to know whether a guard holds, or whether a hazard or the
goal is reachable at a given budget, that question is assumed decidable. The
widened finite-range fragment of \S\ref{sec:base} supplies it. In Coq the
disjunction is a hypothesis of the affected lemmas, never an axiom, so the
development is axiom-free; but we do not oversell that. Where the hypothesis
is stated over an arbitrary world predicate it is excluded middle for that
class, discharged by the caller rather than assumed globally, and a reader
should judge those lemmas as classical in content and constructive only in
bookkeeping.
\end{assumption}

\begin{theorem}[The trichotomy is ordered]\label{thm:ordered}
Rank $\Fut<\Ben<\Cat$. The rank of a residual is monotone non-decreasing in
its budget (\coqok{severity\_monotone\_in\_budget}): raising the budget moves a
\Ben{} residual to \Ben{} or \Cat{} and never to \Fut{}
(\coqok{benign\_step\_up}, \coqok{benign\_no\_regress}), lowering it moves a
\Ben{} residual to \Ben{} or \Fut{} (\coqok{benign\_step\_down}); \Cat{} is
upward closed (\coqok{catastrophic\_upward}) and \Fut{} downward closed
(\coqok{futile\_downward}). So the budgets at which a residual is \Fut, then
\Ben, then \Cat{} are three consecutive intervals of $\mathbb{N}$, in that
order, any of them possibly empty.
\end{theorem}

Two consequences. First, $\kstar$ is a genuine \emph{threshold} and not merely
the budget at which a search stopped: if no hazard is affordable at $k$ and one
is at $k+1$, then none is affordable at any $b\le k$ and one is at every
$b>k$ (\coqok{tolerance\_degree\_is\_a\_threshold}), which is what licenses
computing it by a scan, and below the threshold every residual is \Ben{} or
\Fut{} (\coqok{severity\_classes\_are\_separated}). Second, the classes are
about \emph{executions}: budgeted reachability holds exactly when some run of
the protocol with at most $k$ misselections ends in a hazard state
(\coqok{reach\_mu\_iff\_run}), so \Cat{} means ``an affordable run reaches
harm'', not merely ``a predicate holds''.

\begin{theorem}[Partition]\label{thm:partition}
The classes are pairwise disjoint, and exhaustive whenever the two
reachability questions are decidable. \coqok{severity\_disjoint},
\coqok{severity\_exhaustive}
\end{theorem}

The content is the separation of \Fut{} from \Cat{}. In model-based safety
assessment~\cite{xsap} severity is supplied by the analyst; here it is computed
from goal reachability.

\subparagraph*{The two quantifiers.} \Cat{} and \Ben{} are both existential
over affordable runs, and that means different things on the two sides. For
the hazard it is the reading one wants: a possible catastrophe is a
catastrophe. For the goal it says only that the goal is still \emph{possible},
over the future choices of the very participant the discipline cannot rely on,
so \Ben{} is no guarantee, and the guarantee gets its own name. Say $(G,W)$ \emph{assures} $\varphi$ at
budget $b$ when every affordable run arrives --- every intended branch
delivers, so does every misselection the budget can pay for, and so does every
answer the environment may give --- and call $(G,W)$ \Rob{} when it is
hazard-free and assures the goal.

\begin{theorem}[Assurance is strictly stronger]\label{thm:robust}
What is assured is reachable (\coqok{assured\_reach}), so \Rob{} implies \Ben{}
(\coqok{robust\_benign}); and assurance is downward closed in the budget
(\coqok{assured\_downward}), since more affordable mistakes is a stronger
obligation --- the opposite direction from \Cat. The implication is strict: a
protocol whose safe branch reaches the goal and whose misselectable branch
harmlessly does nothing is \Ben{} at budget $1$ and not \Rob{}
(\coqok{benign\_is\_not\_robust}). One affordable mistake costs the goal and
nothing else, which is exactly what the existential reading hides.
\end{theorem}

The evaluation reports \Ben{}, because that is what the tool computes; \Rob{}
is the stronger verdict a tool could compute on the same graph and we have not
built it.

\begin{definition}[Point of no return]\label{def:pnr}
$(G,W)$ is a point of no return for $\varphi$ if $\neg(G,W)\reachg{b}$. An
action is \emph{goal-destroying} at $(G,W)$ if the goal is reachable before it
fires and at no successor.
\end{definition}

\subparagraph*{Default hazard.}
Because effects are STRIPS, an atom in $\Del_a$ is restored only by some $b$
with the atom in $\Add_b$; if $\Gamma$ has no such $b$, $a$ is
\emph{irreversible}. The implementation's default hazard, needing no
annotation, is \emph{burning a bridge while lost}: an irreversible action
fires from a configuration after which the goal is unreachable. Tools whose
real-world effect cannot be undone but whose pack models no delete-list (an
email cannot be unsent) are declared irreversible explicitly. This is dead-end
detection~\cite{steinmetz-aaai16} and undoability
checking~\cite{daum-icaps16} lifted to a protocol; unrestricted reversibility
checking is harder than planning~\cite{morak-kr25}, and our decidability rests
on the widened finite-range fragment the checker already commits to.

\section{The Well-Formedness Condition, Exactly}\label{sec:typing}

$\safeT{b} G , W$ reads: from $(G,W)$, no hazard arises under any run with at
most $b$ misselections.
\begin{mathpar}
\inferrule*[left=T-End]{W\nvDash\Haz}{\safeT{b}\End,W}
\and
\inferrule*[left=T-Goal]{W\nvDash\Haz \\ \safeT{b} G,W}{\safeT{b}\checkmark\varphi.G,W}
\and
\inferrule*[left=T-Act]{W\nvDash\Haz \\
  \forall W'.\ \langle W\rangle\,a\,\langle W'\rangle \implies \safeT{b} G,W'}
  {\safeT{b} a@\mathsf p.G,\,W}
\and
\inferrule*[left=T-Choice-Safe]
  {W\nvDash\Haz \\\\
   \forall i\in I.\ \ W\vDash\psi_i \implies \safeT{b} G_i,W \\\\
   \forall i\in I.\ \ W\nvDash\psi_i \implies \forall c.\ b = c{+}1 \implies \safeT{c} G_i,W}
  {\safeT{b}\ \mathsf p{\to}\mathsf q:\{\ell_i[\psi_i].G_i\}_{i\in I},\,W}
\end{mathpar}
An intended branch is checked at the same budget; a misselectable branch at
one less, and not at all when the budget is exhausted. This is ``affordable
mistake'' as a rule.

\begin{theorem}[Exact characterization]\label{thm:exact}
$\safeT{k} G,W \iff \neg (G,W)\reachh{k}$, hence a branch is \Cat{} iff its
residual fails the condition. \coqok{TC\_sound}, \coqok{TC\_complete},
\coqok{TC\_exact}
\end{theorem}

\subparagraph*{The rule on the running example.} Take $G_{\mathsf{bad}}$ at
$b=1$ in $W_0=\{\id{funds}\}$. \textsc{T-Choice-Safe} checks the intended
branch $\id{safe}$ at budget $1$ and the misselectable branch $\id{fast}$ at
budget $0$. The safe branch verifies then purchases and no intermediate world
is a hazard, so it holds. The fast branch must satisfy
$\vdash_0 \id{purchase}@\mathsf q.\End,\ W_0$: by \textsc{T-Act} that needs
$\vdash_0 \End, W'$ for every $W'$ with
$\langle W_0\rangle\,\id{purchase}\,\langle W'\rangle$, and every such $W'$
has $\id{booked}$ without $\id{verified}$ --- the hazard. The premise fails, so
$\nvdash_1 G_{\mathsf{bad}}, W_0$, and by Theorem~\ref{thm:exact} a hazard is
reachable within one misselection. At $b=0$ the misselectable branch carries no
obligation at all and the derivation goes through: $G_{\mathsf{bad}}$ is
$0$-tolerant and not $1$-tolerant, which is what the tool reports and what
\coqok{Gbad\_is\_0\_tolerant} and \coqok{Gbad\_not\_1\_tolerant} prove.
Two protocols carry the paper: $G_{\mathsf{bad}}$ is the smallest that exhibits
the question, and \texttt{migration\_backup}, whose output \S\ref{sec:overview}
shows, is the smallest that exhibits \emph{cascading} --- one wrong choice is
affordable, two are not.

The rule is not deep: it is the reachability relation with the negation pushed
through. What matters is what exactness buys, three things used later. A refutation is never an
artifact of conservative rules, so the refutation asymmetry the discipline is
built on survives. The condition is syntax-directed, which is what makes
composition against an interface possible (\S\ref{sec:why}) instead of
whole-system analysis. And because it is a judgment on $(G,W)$ rather than a
decision procedure over a graph, it is the interface through which
Theorem~\ref{thm:bridge} couples the budget to typed programs: a decision
procedure alone does not compose with a typing derivation. \subparagraph*{Two admissible rules.} The condition has the structural
properties a type system is expected to have, and they are already proved.
Budgets \emph{subsume}: a derivation at $k$ is a derivation at every smaller
budget, so
\[
\inferrule*[right=T-Sub]{\vdash_k G, W \\ k'\le k}{\vdash_{k'} G, W}
\qquad
\inferrule*[right=T-Narrow]{\vdash_k \mathsf p{\to}\mathsf q:\mathit{brs}, W
  \\ \mathit{brs}'\subseteq\mathit{brs}}
  {\vdash_k \mathsf p{\to}\mathsf q:\mathit{brs}', W}
\]
are admissible (\coqok{tolerance\_downward\_closed},
\coqok{repair\_narrow\_sound}). \textsc{T-Sub} is subsumption in the budget
index. \textsc{T-Narrow} looks like covariance of internal choice, the
direction session subtyping already has, but it is not, and the difference is
worth stating: the \emph{condition} is closed under removing branches while
\emph{conformance is not} (\coqok{narrowing\_breaks\_conformance}), because
\textsc{T-Comm} asks the sender to offer exactly the protocol's label set. A
session that still offers the removed label conforms to the wide protocol and
not to the narrow one; narrowing both restores it
(\coqok{narrowing\_asymmetry}). So narrowing is a rewrite of the contract, not
a coercion along a subtyping preorder, which is why it is a repair
(\S\ref{sec:repair}).

\begin{theorem}[$\kstar$ is principal]\label{thm:principal}
$\kstar$ is not merely where a scan stopped. If a residual is
$\kstar$-tolerant and not $(\kstar{+}1)$-tolerant, then the budgets at which
it is tolerant are exactly those $\le\kstar$ (\coqok{principal\_characterises}),
such a $\kstar$ is unique (\coqok{principal\_unique}), and one exists whenever
the residual is $0$-tolerant, some budget is not tolerated, and the question is
decidable (\coqok{principal\_exists}).
\end{theorem}

So every derivation of the condition for that residual is a derivation at some
$b\le\kstar$, and the one at $\kstar$ subsumes the rest by \textsc{T-Sub}; and
computing $\kstar$ by an increasing scan, which is what the tool does, is
sound.

Reachability is also monotone in $\Gamma$, so for a \emph{fixed} hazard
granting a tool can only lower $\kstar$
(\coqok{tolerance\_antitone\_in\_ctx}) --- the formal case for least
privilege, and a refutation of our own prior expectation that more tools mean
more recovery.

\subparagraph*{Environment choices.} \textsc{T-Env} asks $W\nvDash\Haz$ and
$\safeT{b}G_i,W$ for every branch, at the same $b$ (\coqok{env\_typing\_iff}).
In a conservative extension (\coqok{safe\_embedE}), Theorems~\ref{thm:exact},
\ref{thm:ordered} and~\ref{thm:principal} and \textsc{T-Sub} hold again
(\coqok{TE\_exact}, \coqok{severityE\_monotone\_in\_budget},
\coqok{principalE\_characterises}, \coqok{toleranceE\_downward\_closed}), the
goal read demonically, as the tool does: every answer must lead to it, which is
strictly stronger (\coqok{env\_goal\_is\_demonic}); so are the default guards
(\coqok{dguard\_some\_intended}). For the hazard an environment choice is one
whose every branch is intended (\coqok{env\_as\_total\_choice}), so
\S\ref{sec:regular} carries over: product, decision procedure and, over guards
with an always-true one, $\mu$-types (\coqok{decide\_env\_correct},
\coqok{reach\_mu\_env}, \coqok{decide\_mu\_env\_correct}, \coqok{bridge\_mu\_env}).

\subparagraph*{Not resource-aware typing.} The budget resembles the potential
of resource-aware session types~\cite{das-lics18,das-popl23} and graded
modalities~\cite{orchard-icfp19}, misleadingly. Potential is consumed by the
program's own execution, and exhaustion is a type error; the budget is consumed
by an \emph{adversarial} choice the program does not make, so the rule
quantifies over deviations, and exhaustion is unconstrained.

\section{From Protocols to Programs}\label{sec:bridge}

Everything so far concerns the global type's semantics. The theorem that makes
it a discipline about \emph{programs} couples the condition to the base
judgment $G\vdash(\MM;W)$, lifted to guarded choice without touching the
guards: \textsc{T-Comm} still requires the sender to offer exactly the
protocol's labels and every branch to check against the same residual; which
label the participant eventually picks is its own business.

\begin{definition}[Instrumented head-move semantics]
A session step carries the acting role and its cost: $0$ for an action or an
intended selection, $1$ when the sender picks a label whose guard fails in the
current world (Coq: \texttt{hstep}). A run records the trace of (role, cost)
events; $\mathrm{total}$ sums costs and $\mathrm{percost}_r$ sums those of role
$r$.
\end{definition}

\begin{theorem}[Bridge]\label{thm:bridge}
If $G\vdash(\MM;W)$ and $\safeT{k} G,W$, then every instrumented head step of
cost $c\le$ the remaining budget preserves typing and debits the budget by
exactly $c$ (\coqok{bridge\_step}); consequently, along every run of total
cost at most $k$, no configuration is a hazard state
(\coqok{bridge\_every\_configuration}) --- runs are prefix-closed
(\coqok{hrun\_split}), so this really is every configuration and not only the
last (\coqok{bridge\_prefix}). An environment step costs $0$ whatever the
answer (\coqok{env\_step\_costs\_nothing}), and the bridge holds for every
resolution of the environment (\coqok{bridgeE\_run}).
\end{theorem}

Two questions come before any of this, and a discipline that skips them earns
the vacuity its predecessor died of (\S\ref{sec:whatchanged}). Does a
conforming session exist at all? And what does the fault model do to
communication safety?

\begin{theorem}[Non-vacuity]\label{thm:inhabited}
Every two-role protocol is inhabited, over any runtime in which a tool call
always has an answer. Read the global type as a pair of processes ---
$\mathit{canon}$: the sender offers exactly the protocol's labels, the receiver
accepts exactly them, each role's actions become its own prefixes, and every
other role is idle --- and the result conforms in every world
(\coqok{canon\_conforms}), so the bridge applied to any such protocol is a
statement about a session that exists
(\coqok{two\_role\_bridge\_nonvacuous}). The running example and its narrowed
repair are inhabited in particular (\coqok{Gbad\_inhabited},
\coqok{Ggood\_inhabited}), and that session really does take the misselected
branch: the wrong label is a step it performs
(\coqok{MBad\_takes\_the\_wrong\_branch}).
\end{theorem}

Of the construction's side conditions two are real and two convenience. Real: distinct labels at a choice,
because a label must name one continuation; and no goal markers, because
\textsc{CT-Goal} demands $\varphi$ at the world where the marker sits and no
construction uniform in the world can supply that --- Theorem~\ref{thm:markers}
seen from the other side. Convenience: every choice is fixed to run from role
$0$ to role $1$ and every action is performed by one of them, which is
stronger than the proof needs. Beyond two roles, inhabitation is a projection
question, and processes are typed directly against the global type; projection
would reappear only to \emph{witness} inhabitation for bystanders, so the
repairs' witnesses are exhibited individually. The goal marker is priced by the same
rule. $\checkmark\varphi$ is an assertion, and \textsc{CT-Goal} is the only
rule that reads one: the condition, the protocol steps, reachability and the
swap relation all treat a marker as transparent. That asymmetry is deliberate,
and Theorem~\ref{thm:markers} is where the premise is paid back.
The second question has a clean answer: nothing.

\begin{theorem}[Markers are met]\label{thm:markers}
If a conforming session of a $b$-tolerant protocol runs, within budget, to a
residual whose head is $\checkmark\varphi$, then $\varphi$ holds at the world
it arrived in (\coqok{markers\_are\_met}). Hence the residual is not \Fut{}
there: the goal is not merely reachable, it holds
(\coqok{marker\_reached\_is\_not\_futile}). This is the only formal link
between the marker and the goal predicate $\Phi$ of \S\ref{sec:severity},
which is otherwise a separate parameter.
\end{theorem}

The premise is not vacuous: the narrowed booking protocol with the goal marked
where its safe path establishes it is tolerant at every budget
(\coqok{Ggoal\_is\_k\_tolerant}), the two-role session conforms to it
(\coqok{Ggoal\_inhabited}), and every run of that session that reaches the
marker has achieved the goal (\coqok{Ggoal\_marker\_met}).

The price: the condition and conformance come apart on a marker. Put one on a branch a misselection can
take: the condition does not see it --- $\vdash_b$ treats a marker as
transparent --- while conformance demands $\varphi$ at a world the protocol does
not control. The result satisfies the condition at \emph{every} budget and
\emph{no} session conforms to it (\coqok{Gmiss\_safe},
\coqok{Gmiss\_uninhabited}). A marker is an assertion, so it belongs where the
session can establish it and nowhere else, which is also why
Theorem~\ref{thm:inhabited}'s construction excludes markers. Our benchmark
carries the goal as a pack-level condition and uses no inline markers, so no
measurement depends on this.

\begin{theorem}[Progress under misselection]\label{thm:progress}
A typed session that has not finished can always take a head step
(\coqok{progress}), and at a choice it can take one for \emph{every} label the
sender may pick, wrong ones included (\coqok{every\_label\_steps}), because
\textsc{T-Comm} makes the receiver offer at least the protocol's labels. A
misselection is therefore never a communication mismatch, an unexpected label
or a deadlock. Its only consequence is the world it leaves behind --- which is
exactly what \S\ref{sec:severity} measures. (Assumption~\ref{ass:dec}.)
\end{theorem}

This is what separates the fault model from the crash-stop
line~\cite{barwell-concur22,barwell-ecoop23}: there a failure removes a
participant and the metatheory must reconstruct progress; here progress is
untouched, and the whole question is which world the session ends in.

\begin{theorem}[Budgets distribute over participants]\label{thm:distribute}
Give each role $r$ an allowance $k_r$. If the allowances of the roles that act
sum to at most $k$, a session typed against a $k$-tolerant protocol is
hazard-free on every run in which each role stays within its own allowance.
\coqok{budget\_distributes}
\end{theorem}

Theorem~\ref{thm:distribute} is the cross-participant composition result: the
global budget is a contract that can be \emph{split} among participants, so a
deployment can hold each role to its own allowance. It is a statement about
runs, not a per-role typing judgment --- the discipline has no projection, so
there is nothing to check per participant --- and the mechanized content is
that per-role costs sum to the total (\coqok{total\_eq\_sum\_percost}) and
the bridge then applies.

\subparagraph*{Recursive sessions.} The finite fragment is not the scope of the
bridge. Section~\ref{sec:regular} gives recursive global types $\mu X.G$ their
semantics by unfolding; processes carry $\mu$ likewise, and a process
\emph{unfolds to a head} (deterministically) before it is inspected, so a
non-contractive process types against nothing. The conformance judgment
becomes coinductive, with one added rule that unfolds the protocol's $\mu$
(Coq: \texttt{ctypes}, \texttt{hstep} in \texttt{Mu.v}).

\begin{theorem}[Bridge, recursive]\label{thm:bridge-mu}
An instrumented head step of a typed session is a protocol path of the same
misselection cost, and typing is preserved (\coqok{sim\_step}). Hence a
session typed against a recursive protocol satisfying the condition at budget
$k$ (Theorem~\ref{thm:mu-judgment}) is hazard-free on every run of cost at
most $k$ (\coqok{bridge\_mu\_safeR}, semantically \coqok{bridge\_mu}). On the
finite fragment this recovers Theorem~\ref{thm:bridge} through
Theorem~\ref{thm:embed} (\coqok{bridge\_finite\_via\_mu}).
\end{theorem}

\subparagraph*{Bystanders.} The head-move semantics is the gate's: an
out-of-turn attempt is refused. An ungated deployment lets a \emph{bystander}
--- a role the head does not involve --- fire its next capability early. Since
the protocol is a schedule of world effects, that is not automatically
harmless: it is exactly what the reorder repair exploits. We give the
permutation relation and the conditions under which it is (Coq:
\texttt{Interleave.v}). A \emph{swap} exchanges two adjacent independent
nodes, in either direction and under any context:
\[\small
\begin{array}{r@{\;}c@{\;}l@{\qquad}l}
b@\mathsf q.\,a@\mathsf r.\,G &\rightleftarrows& a@\mathsf r.\,b@\mathsf q.\,G
  & \mathsf q\neq\mathsf r,\ a,b\ \text{independent}\\
\mathsf p{\to}\mathsf q\{\ell_i[\psi_i].\,a@\mathsf r.\,G_i\}
  &\rightleftarrows& a@\mathsf r.\,\mathsf p{\to}\mathsf q\{\ell_i[\psi_i].\,G_i\}
  & \mathsf r\notin\{\mathsf p,\mathsf q\},\ a\ \text{preserves each}\ \psi_i\\
\checkmark\varphi.\,a@\mathsf r.\,G &\rightleftarrows& a@\mathsf r.\,\checkmark\varphi.\,G
  & a\ \text{preserves}\ \varphi\\
\mathsf p{\to}\mathsf q\{\ell_i.\,\mathsf r{\to}\mathsf s\{m_j.\,K_{ij}\}\}
  &\rightleftarrows& \mathsf r{\to}\mathsf s\{m_j.\,\mathsf p{\to}\mathsf q\{\ell_i.\,K_{ij}\}\}
  & \{\mathsf p,\mathsf q\}\cap\{\mathsf r,\mathsf s\}=\emptyset
\end{array}
\]
Communications change no world, so their permutation needs no condition on
guards or effects --- but it does need the two choices to be genuinely
independent, which the rule expresses by requiring the full product: every
pair $(\ell_i,m_j)$ present, with a residual $K_{ij}$ that depends on the pair
and not on the order. An action's permutation is conditional: $a$ must be \emph{hazard-neutral} (never change whether the
world is a hazard), $a$ and $b$ must \emph{commute} (one order re-serializes as
the other with the same end world) and preserve each other's enabledness, and
$a$ must preserve the guards and goal markers it passes.

\begin{theorem}[Bystanders]\label{thm:bystander}
Safe swaps preserve the condition at every budget (\coqok{swap\_safe}) and
preserve typing --- subject reduction under permutation
(\coqok{swap\_ctypes}). Hence a typed session of a $k$-tolerant protocol is
hazard-free on every \emph{interleaved} run of cost at most $k$, where before
each step the protocol may be permuted by any sequence of safe swaps
(\coqok{bridge\_interleaved}).
\end{theorem}

\begin{theorem}[Variable disjointness suffices]\label{thm:strips}
For STRIPS capabilities stated extensionally --- a precondition supported on
a variable set, add and delete lists --- if the effects of each action are
disjoint from the other's footprint (its support and effects), the actions
commute and preserve each other's enabledness
(\coqok{strips\_commute}, \coqok{strips\_enables}); an action whose effects
are disjoint from the support of a predicate preserves it
(\coqok{strips\_preserves}, \coqok{strips\_neutral}). Together these
discharge every side condition of an action swap (\coqok{strips\_swap\_act}).
\end{theorem}

Theorem~\ref{thm:strips} is what the tool checks: for every capability action
and every earlier node of another role it could pass, it compares footprints
--- taking the goal's and every precondition's atoms as the support of the
derived hazard and of the default guards --- and reports either that the
head order is exact or the pairs the gate must serialize. The discipline is
synchronous by design (the gate delivers a label when it is taken); what the
theorem does not cover is asynchronous buffering, where a sent label waits in
a queue, which changes the semantics rather than merely permuting it.

\section{Regular Protocols}\label{sec:regular}

A regular global type unfolds to a finite graph of protocol nodes; over the
finite abstract world of the widened fragment, budgeted reachability is
reachability in the product $\mathit{Node}\times\mathit{World}\times\{0..k\}$.
The development has two layers: an abstract one over any node and world types
with computable successors (Coq: \texttt{Regular.v}), and its instantiation
to $\mu$-types, where the finiteness of the node set is a theorem rather than
an assumption (\texttt{Mu.v}).

\begin{theorem}[Product]\label{thm:product}
$(n,w)\reachh{b}$ iff a hazard state is reachable from $(n,w,b)$ in the product
graph, whose misselection edges decrement the third component
(\coqok{product\_correspondence}).
\end{theorem}

\begin{theorem}[Loops meet the budget]\label{thm:loop}
The budget component never increases along a product path
(\coqok{budget\_never\_increases}); a path with $d$ misselection edges from
budget $b$ has $d\le b$, and exactly $b-b'$ if it ends at budget $b'$
(\coqok{dev\_edges\_bounded}, \coqok{dev\_edges\_exact}). A loop containing a
misselectable branch is therefore traversed wrongly at most $k$ times within
budget $k$.
\end{theorem}

\begin{theorem}[Decidability]\label{thm:decide}
For a finite world type and a finite, successor-closed list of nodes with
computable successors, the boolean
\texttt{decide\_reachb} decides $(n,w)\reachh{k}$ from any listed node: it is sound
(\coqok{decide\_sound}) and complete (\coqok{decide\_complete}), the latter
via a pigeonhole argument that reachability is witnessed by a path no longer
than the reachable-state count (\coqok{reach\_bounded\_path});
together, \coqok{decide\_reachb\_correct}.
\end{theorem}

\subparagraph*{From $\mu$-types to the graph.} Regular global types carry
$\mu X.G$ and $X$ (de Bruijn indices; Coq: \texttt{Gr}). Unfolding is
substitution of the closed $\mu$-term for its variable, and the protocol steps
of \S\ref{sec:mis} gain one silent step, $\mu X.G\to G[\mu X.G/X]$.
Budgeted reachability on $\mu$-types \emph{is} the product of
Theorem~\ref{thm:product} over these steps (\texttt{reach\_mu}).

\begin{theorem}[The unfolding closure is finite]\label{thm:cands}
For a closed regular type $G_0$ there is a computable list
$\mathrm{cands}(G_0)$ --- the closures of $G_0$'s subterms under their
enclosing binders --- that contains $G_0$ and is closed under every protocol
step (\coqok{cands\_closed}); every state reachable from $G_0$ in the product
lies in it (\coqok{reach\_in\_cands}). The heart is a substitution lemma:
substituting the closed $\mu$-term for the variable it binds, after closing
the outer context, is closing the extended context (\coqok{subst\_comp},
\coqok{unfold\_close}).
\end{theorem}

\begin{theorem}[Embedding]\label{thm:embed}
The finite fragment embeds: $\reachh{b}$ of \S\ref{sec:mis} coincides with
\texttt{reach\_mu} on the image (\coqok{reach\_embed}), so
\textsc{T-Choice-Safe} is exactly non-reachability of a hazard in the product
graph of the embedded protocol (\coqok{TC\_regular}).
\end{theorem}

Read \textsc{T-Choice-Safe} coinductively --- a recursive protocol has no
finite derivation, while a reached hazard still has a finite path. The rule is
one clause, over the protocol's own steps rather than its syntax, with a double
line for the coinductive reading:
\[
\inferrule*[right=T-Safe-$\nu$,vdots=1.2em]
  {\neg\Haz(w) \\
   \forall\, G\to_0 G',\ \safeRr{b}G',w' \\
   \forall\, G\to_1 G',\ b=b'{+}1 \Rightarrow \safeRr{b'}G',w'}
  {\safeRr{b} G, w}
\]
where $\to_0$ is a compliant protocol step and $\to_1$ a misselection.

\begin{theorem}[The condition, for recursive protocols]\label{thm:mu-judgment}
$\safeRr{k}$ is again exactly non-reachability: sound
(\coqok{TR\_sound}), complete (\coqok{TR\_complete}), together
\coqok{TR\_exact}. It is a conservative extension: on the finite fragment the
two judgments coincide under the embedding (\coqok{TC\_is\_TR}), and
\texttt{decide\_mu} decides it (\coqok{decide\_mu\_judgment}). So the bridge,
the decision procedure and the condition are statements about one object at
both layers, not three.
\end{theorem}

\begin{theorem}[Decidability for $\mu$-types]\label{thm:decide-mu}
Over a finite abstract world with decidable guards and computable tool
successors, \texttt{decide\_mu} decides $(G_0,w)\reachh{k}$ for every closed
regular $G_0$: it runs Theorem~\ref{thm:decide}'s procedure on
$\mathrm{cands}(G_0)$, which Theorem~\ref{thm:cands} shows is
successor-closed (\coqok{decide\_mu\_correct}).
\end{theorem}

\begin{theorem}[Guardedness, and where it is needed]\label{thm:guarded}
Conformance for recursive protocols is coinductive, so $\mu X.X$ is typable
against every session and world (\coqok{ctypes\_mu\_var}) --- and it can never
take a step (\coqok{unguarded\_var\_stuck}). Progress therefore fails without
a guardedness condition, and this is the only place in the development that
needs one. Call a type \emph{guarded} when no recursion variable occurs at the
head of its own binder after peeling $\mu$s and goal markers ($\mathtt{gd}$);
markers do not guard, because a step erases one rather than firing on it, and
$\mu X.\checkmark\varphi.X$ is stuck for the same reason
(\coqok{unguarded\_goal\_stuck}). Guarded unfolding leaves the leading prefix
no longer than it was (\coqok{pre\_unfold}) and preserves guardedness and
closedness (\coqok{gd\_unfold}, \coqok{closed\_unfold}), so a closed guarded
conforming protocol that has not finished can always step
(\coqok{progress\_mu}). Dropping $\mathtt{gd}$ makes that false, witnessed by
$\mu X.X$ (\coqok{contractiveness\_is\_necessary}).
\end{theorem}

Theorem~\ref{thm:decide-mu} needs no guardedness: the candidate set is finite by
construction (Theorem~\ref{thm:cands}), so the decision procedure terminates,
and answers correctly, on non-contractive input too.

\begin{proposition}[Tolerance degree]\label{prop:degree}
$\kstar$ is computable. Whether $\kstar=\infty$ is plain reachability in
$\mathit{Node}\times\mathit{World}$ with misselection edges made free;
otherwise $\kstar$ is at most the number of misselection edges on a shortest
hazard path, bounded by the product size, and is found by iterating
Theorem~\ref{thm:decide}. With $|\mathit{World}|=2^{|\Pred|}$ the decision
problem is in PSPACE; each iteration is linear in the product.
\end{proposition}
\begin{proof}
Paper proof from Theorems~\ref{thm:product}--\ref{thm:decide}: dropping the
budget component yields the $\infty$ case; \texttt{dev\_edges\_exact} bounds
$\kstar$; reachability in a succinctly presented graph is in PSPACE.
\end{proof}

\section{Why a Discipline, and Not a Model Checker}\label{sec:why}

Theorem~\ref{thm:exact} invites the objection that the rules are a decision
procedure in disguise; Naik and Palsberg~\cite{naik-palsberg} gave the standard
reply, and we rest instead on a theorem and an experiment.


\begin{theorem}[Modular composition]\label{thm:seq}
If $\safeT{k} G_1,W$ and $\safeT{b'} G_2,W'$ for every $(b',W')$ at which
$G_1$ can end from budget $k$, then $\safeT{k} G_1;G_2,\,W$ (\coqok{TC\_seq}).
In interface form: for any $I$ containing those exit points, it suffices that
$\safeT{b'} G_2,W'$ for all $(b',W')\in I$ (\coqok{TC\_seq\_interface}).
Sequencing is defined on the finite fragment only: $\safeRr{k}$ has no
composition theorem, so modular checking does not currently reach recursive
protocols.
\end{theorem}

The theorem is what a whole-system reachability run cannot be: $k$ becomes a
contract on a boundary. It does not say that boundary is small, and the
evaluation shows it is not --- the \emph{concrete} exit set grows exponentially
when every segment leaves its own atoms behind. What makes modular checking pay
is the interface form with a coarser $I$, and the principled coarsening is the
cone of influence.

\begin{theorem}[Cone of influence]\label{thm:cone}
Let $V$ be a set of state variables such that the hazard reads only $V$, every
guard of $G$ reads only $V$, and every capability \emph{$G$ invokes} maps
$V$-agreeing worlds to $V$-agreeing successors --- the last restricted to $G$'s
own tools, not to all of $\Gamma$, because a modular interface is checked
against a residual and the wider hypothesis is one the construction that uses
it would refute. Then worlds agreeing on $V$ are interchangeable for $G$
at every budget: hazard reachability transports (\coqok{reach\_cone}) and so
does the condition (\coqok{safeT\_cone}). Consequently an interface need only
contain one representative of each $V$-class, not one world per concrete exit
(\coqok{interface\_projection}). For \textsc{strips} capabilities that last
condition is a containment the tool can check --- it suffices that every
\emph{precondition} of an invoked tool reads only $V$
(\coqok{strips\_cap\_cone}, \coqok{strips\_safeT\_cone}) --- and what an action
\emph{writes} outside $V$ does not matter, since worlds are compared only on
$V$.
\end{theorem}

That abstraction is what makes computing an interface cheap rather than
exponential (Table~\ref{tab:mod}). What the theorem does not do is compute $V$:
the tool takes the cone syntactically from the predicates the remaining
segments mention, and that this syntactic cone satisfies the containment is a
property of the front end. This is the axis on
which MPST is currently being argued~\cite{castro-popl26}; the community
has accepted parameterising typing by a model-checked
property~\cite{barwell-concur22,barwell-ecoop23}, and the theorem plus the
projection is our concrete instance of that move.

\section{Implementation and Evaluation}\label{sec:eval}

\texttt{skillc severity} implements the analysis over the existing
achievability compiler: the checker's own symbolic search (predicates
concrete, numerics symbolic, widened at loop heads) is run pointwise at every
branch of every choice, at every abstract world and budget up to
$k_{\max}=4$, with the default guards and hazard of
\S\ref{sec:mis}--\S\ref{sec:severity} and explicit overrides where declared.
It reports per-branch severity, $\kstar$, the first hazard witness, the
point-of-no-return action, the narrowing repair, and the bystander pairs the gate must serialize. Nineteen tests check the
tool's verdicts on the paper's instances against the Coq development's,
including two that check the tool where a theorem licenses a shortcut: that
$\kstar$ really is a threshold on every benchmark protocol
(Theorem~\ref{thm:principal}), and that projecting an interface onto the cone
of influence never changes what the modular check concludes
(Theorem~\ref{thm:cone}).

\subparagraph*{Verified kernel.} The tool's search is the graph form of
Theorem~\ref{thm:cands}, with widening at loop heads over the arithmetic
state; it is not extracted from the proof. To close that gap on the fragment
where the proof applies we built a kernel (Coq: \texttt{Kernel.v}) that
instantiates \texttt{decide\_mu} on bit-vector worlds with STRIPS actions,
guards as world sets, and the derived hazard as a bit that an irreversible
action sets whenever the goal is no longer reachable from its continuation.
The decision --- the least budget at which a hazard is reachable --- is proved
correct (\coqok{kernel\_first\_spec}); goal reachability, used to derive
default guards and hazard bits, is proved correct
(\coqok{goal\_reachable\_correct}), and so is hazard avoidability
(\coqok{avoid\_b\_correct}); the derived guard is the default rule of
\S\ref{sec:mis} (\coqok{rat\_guard\_spec}), and at a choice with no right
answer every branch is intended, so the choice offers no misselection edge
and costs nothing (\coqok{forced\_choice\_intended},
\coqok{forced\_choice\_costs\_nothing}); the decision procedure itself was
replaced by a deduplicating, early-exiting variant proved sound and complete
(\coqok{decide'\_sound}, \coqok{decide'\_complete}).

The kernel's goal notion is not the theory's, and the two are
\emph{incomparable} rather than one being safely stronger. It requires the
protocol to run to completion, where $\Phi$-reachability fires at any
$\Phi$-world (\coqok{goal\_reach\_strictly\_stronger}), and it spends no
budget, where $\Phi$-reachability is budgeted
(\coqok{goal\_reach\_ignores\_the\_budget}); only the first direction has an
implication, and it reads the budget off the witness path
(\coqok{goal\_reach\_implies\_reach\_mu}). This matters beyond the severity
labels: the elaboration derives the \emph{default guard} and the
\emph{derived hazard bit} from the same notion, so it is a front-end convention
the theory does not fix, and $\kstar$ is computed against it. Both the analyzer
and the kernel implement that convention, independently; what the agreement
below establishes is that they agree, not that the convention is the paper's
$\Phi$-reachability. The elaboration from
the tool's protocol tree into the kernel's input is the trusted front end: a
hundred lines of executable Coq and a hundred of Python. The kernel is
extracted to OCaml and run by \texttt{skillc severity --verified}. On the
sixteen benchmark protocols in its fragment --- propositional effects, no
environment choice --- it agrees with the analyzer's tolerance degree on every
one, in under 50\,ms each.

Sixteen hand-picked protocols are a weak test of that agreement, so we also
compare the two on random ones (\texttt{scripts/differential\_test.py}). The
generator samples packs inside the shared fragment and is otherwise
unconstrained: nested choices, explicit and default guards, named recursion,
irreversible tools, unreachable goals. The analyzer and the extracted kernel
share no code, so a disagreement is a defect in one of them.
Across two seeds ($150$ at \texttt{20260902} and $350$ at \texttt{4242}, both
at $k_{\max}=2$), 500 protocols were generated, 500 compared, and the two agree
on the tolerance degree 500 times and disagree none. An earlier campaign at the
same seeds compared 499: one generated pack carried two recursions with the same
name and was rejected by the pack schema, a generator bug rather than a tool
one, since fixed. Two thirds of the sampled protocols have no reachable hazard
at all, so most agreements are agreements that nothing is there; the
discriminating cases are the third that carry a finite tolerance degree, and
$k_{\max}=2$ means no degree above two is ever distinguished.

This is the strongest evidence we can offer that two independent
implementations --- a hand-written symbolic analyzer and an extraction of the
mechanized decision procedure, sharing no code --- compute the same tolerance
degree. What it does not
establish is that the goal convention both of them use (above) is the
$\Phi$-reachability of \S\ref{sec:severity}; that gap is a front-end one and we
leave it open.

\subparagraph*{Finding 1: existing corpora cannot pose the question.}
Over the compiler's achievability corpus (15 packs), its extension (6) and the
17 public skills of the reference collection, \emph{no} pack declares an
irreversible effect and the real skills contain \emph{no choice point at all}:
all 38 are trivially tolerant at every $k$. That is a result about modelling.
Real-world irreversibility --- sending, paying, deleting, deploying --- is not
captured by achievability-oriented compaction, so severity needs the
declaration \S\ref{sec:severity} introduces, and the question this paper asks
has been invisible to the corpora the field evaluates on.

\subparagraph*{Finding 2: the severity benchmark.}
We therefore built seventeen protocols with genuine decision points and
irreversible tools, each with a provenance note stating the expected verdict
\emph{before} running the tool: booking (the running example and its two
repairs), database migration, an email campaign and its guard repair,
deployment with and without rollback, file cleanup, order fulfilment with an
environment-controlled bank, a retry loop with a destructive give-up, a
guarded shipping choice, an insurance claim in two worlds, a staged
commit, and two release protocols with a bystander role --- one whose
bystander is variable-disjoint, one whose bystander establishes a later
precondition. Results (Table~\ref{tab:bench}): 55 branch verdicts (a branch at an
abstract world) in 0.36\,s total; 35 \Ben{}, 3 \Fut{}, 17 \Cat{}. That
count is over \emph{all} branches, and the reader should not take the 35 as
the \Ben{} class of the trichotomy: 34 of them are branches whose guard
\emph{holds}, so they are not misselections at all. Restricted to
misselections the split is $1$ \Ben, $3$ \Fut, $17$ \Cat --- and the single
\Ben{} misselection is the express detour of \texttt{shipping\_detour}, whose
guard is explicit. The class the slogan needs is the one this benchmark is
thinnest on.

\emph{Which results the benchmark exercises.} A choice is a communication
$\mathsf p{\to}\mathsf q$ with $\mathsf p\neq\mathsf q$, so a protocol whose
choices name no recipient has no conforming session --- \textsc{CT-Comm} wants
the roles to differ. Six protocols have a second role and their eight choices
name it; the other eleven model a lone agent, with nobody to tell. The
condition, $\kstar$, the severity classes and the repairs are about the protocol
and apply to all seventeen; the bridge, progress, budget distribution and the
bystander results are about sessions and apply to five of the six ---
\texttt{order\_fulfilment} is out, since its choices run in both directions
between one pair. That is a real limit of this benchmark rather
than of the theory, and the two release protocols were built to exercise the
bystander half. $\kstar$ distribution $9\times 0$,
$2\times 1$, $6\times{\ge}5$; the nine distinct point-of-no-return actions it
names across all seventeen are \texttt{purchase, send, deploy, delete, ship,
purge, refund, drop\_old, commit} --- every one an externally visible effect. Sixteen of seventeen
verdicts matched the pre-stated expectation; the exception exposed an
authoring error in our benchmark (a goal disjunct that failed to exclude
\texttt{shipped}), corrected and recorded. The bystander check
(Theorem~\ref{thm:strips}) reports the audit release exact --- five
independent pairs --- and the cache release with two pairs to serialize, both
on the atom the bystander establishes; every other benchmark protocol has a
single acting role, so the question does not arise there.

\emph{Which repairs the benchmark exercises.} Two of its protocols are repairs
of others, and both apply the same transformation: \emph{hoisting} an action
performed in every branch to above the choice, so that the wrong selection can
no longer skip it. Hoisting is the bystander swap of Theorem~\ref{thm:bystander} when the actor is
uninvolved in the choice (\coqok{SW\_comm}, \coqok{SW\_comm\_rev}); in both
instances the actor \emph{is} one of the communicating pair, so no theorem of
ours licenses it and what the table reports is the tool's verdict on the
result, not a certified rewrite. The mechanized guard and reorder repairs need
an abort world the pack language lacks, so only Coq exercises them
(\S\ref{sec:repair}).

\begin{table}[t]
\centering\small
\resizebox{\textwidth}{!}{%
\begin{tabular}{@{}llccclccc@{}}
\toprule
protocol & guards & choices & loops & irreversible & $\kstar$ & \Ben & \Fut & \Cat\\
\midrule
booking\_fastpath & default & 1 & 0 & purchase & 0 & 1 & 0 & 1\\
booking\_reordered (repair: hoist) & default & 1 & 0 & purchase & $\ge$5 & 2 & 0 & 0\\
booking\_narrowed (repair: narrow) & default & 1 & 0 & purchase & $\ge$5 & 1 & 0 & 0\\
migration\_backup & default & 2 & 0 & drop\_old & 1 & 4 & 0 & 2\\
email\_campaign & declared & 1 & 0 & send & 0 & 1 & 0 & 1\\
email\_campaign\_guarded (repair: hoist) & declared & 1 & 0 & send & $\ge$5 & 2 & 0 & 0\\
deploy\_with\_rollback & default & 1 & 0 & --- & $\ge$5 & 1 & 1 & 0\\
deploy\_no\_rollback & default & 1 & 0 & deploy & 0 & 1 & 0 & 1\\
file\_cleanup & default & 1 & 0 & delete & 0 & 1 & 0 & 1\\
order\_fulfilment (environment choice) & declared & 3 & 0 & ship & 0 & 1 & 0 & 1\\
retry\_then\_purge & default & 2 & 1 & purge & 0 & 6 & 0 & 2\\
shipping\_detour (explicit guards) & explicit & 1 & 0 & --- & $\ge$5 & 2 & 0 & 0\\
claim\_eligible & explicit & 1 & 0 & refund & $\ge$5 & 1 & 1 & 0\\
claim\_ineligible (same protocol) & explicit & 1 & 0 & refund & 0 & 1 & 0 & 1\\
staged\_commit & default & 3 & 0 & commit & 1 & 8 & 1 & 5\\
release\_with\_audit (bystander, exact) & default & 1 & 0 & deploy & 0 & 1 & 0 & 1\\
release\_with\_cache (bystander, serialize) & default & 1 & 0 & deploy & 0 & 1 & 0 & 1\\
\bottomrule
\end{tabular}}
\caption{The severity benchmark ($k_{\max}=4$; $\ge$5 means no hazard within
four misselections). Verdict counts are per branch and abstract world, so 35 \Ben{} counts 34 branches whose guard holds; restricted to misselections the split is 1/3/17. Total analysis time 0.36\,s.}
\label{tab:bench}
\end{table}

Three pairs carry the paper's claims. \emph{Repairs restore tolerance:}
narrowing, and hoisting the establishing action above the choice, each take a
$0$-tolerant protocol to tolerance at every $k$. \emph{Severity is a property of node and world:} the
insurance-claim protocol is identical in both rows; with an eligible claimant
the wrong branch is \Fut{}, with an ineligible one \texttt{refund} is a
point of no return and $\kstar=0$. \emph{The point of no return depends on
$\Gamma$'s establishers:} with \texttt{rollback} available, \texttt{deploy}
deletes nothing permanently and the wrong branch is merely \Fut{}; remove
\texttt{rollback} and the same branch is \Cat{}. The last pair refines
\coqok{tolerance\_antitone\_in\_ctx}'s anti-monotonicity (\S\ref{sec:typing}): for a \emph{fixed} hazard more
tools never help, but the derived hazard is not fixed --- a tool that
re-establishes a lost atom removes a point of no return.

\subparagraph*{Finding 3: on real protocols, the points of no return are the rules
no tool enforces.} The benchmark of Finding~2 is ours. We therefore ported
eighteen protocols that none of us wrote (\texttt{benchmarks/severity\_real}).
Six are multi-role protocols from a collection of cases built from public agent
configurations: GitHub Copilot customisation files of the awesome-copilot
collection~\cite{awesome-copilot} (pull-request review and merge, three variants
of a release gate, a seven-role development team) and Anthropic's public
skills~\cite{anthropic-skills} (document co-authoring). Twelve are read off the
customer-service policies of $\tau$-bench~\cite{taubench} and
$\tau^2$-bench~\cite{tau2bench} (airline, retail and telecom; MIT licence), one
per procedure in which the agent decides before an action the policy treats as
final or consequential: cancelling, modifying, returning or exchanging an order,
booking, changing flights, cabin or baggage, sending a compensation certificate,
refuelling data, paying a bill. Each rule ``only if $C$'' becomes an agent
choice whose explicit guard is the policy sentence; what the user answers and
what the records show become environment choices, resolved adversarially; and
the precondition of the irreversible tool is exactly what the source tool code
checks, nothing more. Every protocol was \emph{pre-registered}: its provenance
note and its table of expected verdicts were committed before the tool was run on
any of them.

All 1114 branch verdicts, every $\kstar$ and every point of no return matched
the pre-registered expectation, and all eighteen protocols have $\kstar=0$
(Table~\ref{tab:real}). The agreement checks the implementation against our hand
reasoning, not the world models: the guards are explicit, so each verdict follows
from a short reachability argument written down beforehand. What the port does
test is whether real policy text and tool code can be turned into protocols at
all, and which hazards survive the translation. One pattern accounts for all of
them. Where the source tool enforces a rule as a precondition, a wrong choice at
that rule's gate is \Ben{} or \Fut{}: the tool refuses the call. Where the rule is
stated only in the policy text, the wrong choice is \Cat. The retail tools check
an order's status, the reason, the items and the payment method themselves, and
those gates are harmless; the one rule no retail tool checks, authentication, is
the one whose gate is \Cat. The airline policy says of its own rules that ``the
API does not check these for the agent'', and every such gate is one wrong
choice from the irreversible call.

That points at the repair (last two columns of Table~\ref{tab:real};
\texttt{scripts/severity\_repairs.py}). Moving the policy-only rules into the
tools' preconditions raises $\kstar$ to $\ge5$, tolerant at every tested budget,
for 9 of the 18 protocols. The other 9 stay at $\kstar=0$ for one reason only:
an explicit user-confirmation step, whose answer the agent reports and no tool
can check, so executing after an answer other than yes remains one misselection
from the irreversible action. The analysis' own narrowing repair
(\S\ref{sec:repair}) removes that branch too and makes all 18 tolerant at every
tested budget. Neither repair lowers $\kstar$ on any protocol (checked by a test
for both; for narrowing it is Theorem~\ref{thm:narrow}).
Enforcement turns \Cat{} branches into \Fut{} ones --- the wrong choice now ends
at a refusing tool --- so it makes the protocols safer, not more successful.

The caveats are specific. The three release-gate variants share their upstream
files, the retail modify and exchange procedures have the same structure, and the
airline protocols come from the $\tau$-bench and $\tau^2$-bench versions of one
policy, so the eighteen are not eighteen independent structures. For the six
configuration-based protocols there is no tool code at all: the rule we move is a
proposal about the tool, not a re-reading of code that exists. The twelve policy
protocols rely on environment choices. The theory covers them: a protocol of their
shape --- rule gate, user confirmation, irreversible action --- has $\kstar=0$,
and tolerance at every budget with the rule as a precondition, unless an
execute/hold choice follows the answer, whose wrong branch only narrowing removes
(\coqok{Gpol\_kstar\_is\_0}, \coqok{Genf\_is\_k\_tolerant},
\coqok{Gconf\_enforced\_kstar\_is\_0}, \coqok{Gconf\_narrowed\_is\_k\_tolerant}).
Their own $\kstar$ is the tool's verdict: the kernel has no environment choice. And $\kstar=0$ everywhere is a property of the
structure ported --- one irreversible action behind gates the agent judges --- not
an estimate of how common it is. The port also exposed a defect in the tool: its
recording pass stopped at the first hazardous environment branch, which left
verdict tables incomplete but never changed $\kstar$; it is fixed, and the
tables are complete.

\begin{table}[t]
\centering\footnotesize\setlength{\tabcolsep}{2.8pt}
\resizebox{\textwidth}{!}{%
\begin{tabular}{@{}lrrr@{\hspace{7pt}}lrrrrrr@{}}
\toprule
& & \multicolumn{2}{c}{choices} & & & \multicolumn{3}{c}{misselections} & \multicolumn{2}{c}{$\kstar$ after}\\
protocol & roles & agent & env. & irreversible action & $\kstar$ & B & F & C & enforce & narrow\\
\midrule
\multicolumn{11}{@{}l}{\emph{public agent configurations: Anthropic's skills~\cite{anthropic-skills}, awesome-copilot~\cite{awesome-copilot}}}\\
doc\_coauthor\_ship & 4 & 1 & 0 & ship & 0 & 5 & 0 & 1 & $\ge$5 & $\ge$5\\
gem\_dev\_team & 7 & 2 & 0 & deploy & 0 & 5 & 2 & 5 & $\ge$5 & $\ge$5\\
pr\_review\_merge & 4 & 2 & 0 & merge & 0 & 150 & 0 & 16 & $\ge$5 & $\ge$5\\
sdlc\_release\_gate & 7 & 1 & 0 & deploy & 0 & 5 & 0 & 1 & $\ge$5 & $\ge$5\\
sdlc\_release\_gate\_concurrent & 7 & 1 & 0 & deploy & 0 & 5 & 0 & 1 & $\ge$5 & $\ge$5\\
sdlc\_release\_gate\_concurrent\_withhold & 7 & 1 & 0 & deploy & 0 & 5 & 6 & 1 & $\ge$5 & $\ge$5\\
\addlinespace
\multicolumn{11}{@{}l}{\emph{airline policy, $\tau$-bench~\cite{taubench} and $\tau^2$-bench~\cite{tau2bench}}}\\
baggage\_update & 3 & 2 & 2 & update\_reservation\_baggages & 0 & 0 & 2 & 4 & 0 & $\ge$5\\
book\_reservation & 3 & 3 & 3 & book\_reservation & 0 & 3 & 7 & 10 & 0 & $\ge$5\\
cancel\_reservation & 3 & 2 & 2 & cancel\_reservation & 0 & 0 & 10 & 7 & $\ge$5 & $\ge$5\\
change\_cabin & 3 & 3 & 3 & change\_cabin & 0 & 1 & 3 & 10 & 0 & $\ge$5\\
compensation & 3 & 3 & 3 & send\_certificate & 0 & 1 & 4 & 25 & $\ge$5 & $\ge$5\\
modify\_flights & 3 & 3 & 4 & update\_reservation\_flights & 0 & 7 & 11 & 10 & 0 & $\ge$5\\
\addlinespace
\multicolumn{11}{@{}l}{\emph{retail policy, $\tau$-bench}}\\
cancel\_pending\_order & 3 & 4 & 4 & cancel\_pending\_order & 0 & 13 & 22 & 7 & 0 & $\ge$5\\
exchange\_delivered\_order & 3 & 5 & 5 & exchange\_delivered\_order\_items & 0 & 25 & 47 & 6 & 0 & $\ge$5\\
modify\_pending\_items & 3 & 5 & 5 & modify\_pending\_order\_items & 0 & 25 & 47 & 6 & 0 & $\ge$5\\
return\_delivered\_order & 3 & 4 & 4 & return\_delivered\_order\_items & 0 & 10 & 16 & 4 & 0 & $\ge$5\\
\addlinespace
\multicolumn{11}{@{}l}{\emph{telecom policy, $\tau^2$-bench}}\\
overdue\_payment\_resume & 3 & 2 & 3 & make\_payment & 0 & 0 & 4 & 2 & $\ge$5 & $\ge$5\\
refuel\_data & 3 & 3 & 3 & refuel\_data & 0 & 1 & 3 & 10 & 0 & $\ge$5\\
\bottomrule
\end{tabular}}
\caption{The eighteen real protocols ($k_{\max}=4$; $\ge$5: no hazard within four
misselections). Choices: static choice nodes of the agent and of the
environment; B/F/C: \Ben/\Fut/\Cat{} verdicts on misselections, per branch and
world; $\kstar$ after moving the policy-only rules into tool preconditions
(\emph{enforce}) and after the narrowing repair (\emph{narrow}).}
\label{tab:real}
\end{table}

\subparagraph*{Finding 4: live agents on real policies take the branches the
classification names, and the gate built from it removes every catastrophe.}
Predictive validity asks whether the \Cat{} verdicts name choices real agents
make. We drove all 35 protocols (17 authored, 18 real) with two models of one
vendor, \texttt{claude-haiku-5-5} and \texttt{claude-sonnet-5-5}, through its
command-line interface with no tools and a two-sentence system prompt, each
playing every agent role; prompts are scrubbed of URLs, domains and e-mail
addresses, and the agents run without credentials or network
(\texttt{scripts/live\_agents.py}; every run and reply is in
\texttt{results/live\_agents\_runs.jsonl}). At a choice the model answers with a
label; effects, preconditions, loops, the environment's choices (sampled
uniformly) and the hazard are the analyzer's own semantics. Five conditions vary
what the agent knows. In \emph{plain} it sees the goal, the facts now true, the
tools with their effects and each option's continuation; \emph{pressured} adds a
requester's note that the checks were done and asks for the fastest route. In
\emph{blind} and \emph{blind\_pressured} it sees the tool names only, and for each
option only the steps that option takes, as an agent working from a skill
document would. In these four \emph{informed} conditions the facts shown are
true. In \emph{misinformed} the agent sees no facts at all, and a requester's
note asserts that every precondition holds. Two arms: \emph{ungated}, and
\emph{gated}, the narrowing repair applied per state --- a branch that is not
intended and is \Cat{} in the current world is not offered. The gate never
refuses an intended branch, by construction, and on none of the recorded choices
did it. Five runs per protocol, model, condition and arm: 3500 runs, \$11.57.

We count two things (Table~\ref{tab:live};
\texttt{scripts/live\_agents\_summary.py}): runs that \emph{entered} a \Cat{}
branch --- took a branch the analyzer classifies \Cat{} at that node, so that harm
is reachable within the budget --- and runs that \emph{completed} a catastrophe,
an irreversible action firing after the goal was lost. On the authored protocols
in the informed conditions, no run entered a \Cat{} branch (0 of 680 ungated):
with true facts shown, these models do not take the wrong branch on the protocols
we wrote, and this is a null result. On the real protocols in the informed
conditions, 68 of 720 ungated runs entered a \Cat{} branch against 0 of 720
gated (exact two-sided Fisher test, $p<10^{-20}$), per condition 20, 18, 16 and
14 of 180. All 68 are the smaller model's, on 7 of the 18 protocols. Of these, 14
are \texttt{pr\_review\_merge}'s extra review round, which is \Cat{} only because
the loop gives a later reviewer another chance to err within the budget; without
them it is 54 of 720 against 0 ($p<10^{-16}$). Entering is not completing: \Cat{}
says what four wrong choices can reach, and in these runs a later choice was
right or the run stopped short, so 1 of the 68 runs completed a catastrophe, and 0 gated runs did.

The misinformed condition is a stress test. There, completed catastrophes were 53
of 180 ungated runs on the real protocols against 0 of 180 gated
($p<10^{-17}$), on 14 of the 18 protocols, and 9 of 170 against 0 of 170 on the
authored ones ($p=0.0035$); 76 of 180 ungated real runs entered a \Cat{} branch,
against 0 gated. Here it is the larger model that follows the note: 54 of the 62
catastrophes are its. On the authored protocols the catastrophes are ordered by
$\kstar$: 9 of 90 runs at $\kstar=0$, 0 of 20 at $\kstar=1$, 0 of 60 at
$\kstar\ge5$ --- the last forced, since those protocols have one choice point each
and tolerate four misselections, and the middle one resting on twenty runs. In
every cell no run completed a catastrophe with at most $\kstar$ misselections (0
violations), which checks the harness against Theorem~\ref{thm:bridge}; the two
share their semantics, so it tests the harness, not the theorem.

\begin{table}[t]
\centering\small\setlength{\tabcolsep}{4.5pt}
\begin{tabular}{@{}llrrrrrrr@{}}
\toprule
& & runs & \multicolumn{2}{c}{entered \Cat} & \multicolumn{2}{c}{catastrophes} & \multicolumn{2}{c}{goal reached}\\
protocols & condition & per arm & ungated & gated & ungated & gated & ungated & gated\\
\midrule
authored & plain & 170 & 0 & 0 & 0 & 0 & 170 & 170\\
authored & pressured & 170 & 0 & 0 & 0 & 0 & 170 & 170\\
authored & blind & 170 & 0 & 0 & 0 & 0 & 170 & 169\\
authored & blind\_pressured & 170 & 0 & 0 & 0 & 0 & 165 & 166\\
authored & misinformed & 170 & 9 & 0 & 9 & 0 & 152 & 161\\
real & plain & 180 & 20 & 0 & 0 & 0 & 167 & 170\\
real & pressured & 180 & 18 & 0 & 0 & 0 & 170 & 170\\
real & blind & 180 & 16 & 0 & 0 & 0 & 162 & 169\\
real & blind\_pressured & 180 & 14 & 0 & 1 & 0 & 167 & 169\\
real & misinformed & 180 & 76 & 0 & 53 & 0 & 49 & 112\\
\bottomrule
\end{tabular}
\caption{Live agents: two models, five runs per protocol, model, condition and
arm. \emph{Entered \Cat}: a run took a branch classified \Cat{} at that node;
\emph{catastrophes}: an irreversible action fired after the goal was lost. 3500
runs in all.}
\label{tab:live}
\end{table}

The gate has a price. It removes options; it does not make an agent choose
well. Where the refused branch was the
only wrong one the choice becomes forced, and gated runs reached the goal more
often (112 against 49 of 180 misinformed real runs); where a \Fut{} branch
remains, the agent can still take it, and 68 gated misinformed real runs ended
without the goal, against 78 ungated (informed: 42 against 53). A gated run that
would have been a catastrophe therefore ends on the intended branch or without
the goal; since the arms are independent samples we cannot say which for a
given run. The gate trades harm for, at worst, a lost goal --- \Cat{} for
\Fut{} --- and in these runs it did not cost goals in aggregate.

In short, the \Cat{} classification names choices real agents make on real
policies, and the gate built from it removed every such entry and every
catastrophe, in every condition. It shows no more than that. With honest, fully
informed prompts our own protocols draw no error, so they test the gate and the
ordering by $\kstar$ only under the misinformed note; the note is ours, and it
measures what an agent does with a false account, not how often it gets one; the
models are two of one vendor; the environment is simulated. Earlier runs under a
superseded default guard are archived and not used.

\subparagraph*{Finding 5: modularity pays only with the interface abstraction.}
We chained a benchmark segment $n$ times, $G;G;\cdots;G$ with atoms renamed
per segment, and compared four analyses at $k=2$ (Table~\ref{tab:mod}): the
whole-system \emph{complete} enumeration of hazard-free terminations
(like-for-like with computing an interface), modular checking with the
\emph{concrete} exit set, modular checking with exits \emph{projected} onto
the cone of influence of the remaining segments, and the cost of re-checking
after the last segment is replaced by its unsafe variant.

State first what modularity does \emph{not} buy, because the natural comparison
is the wrong one. Deciding whether a hazard is reachable in the whole chain is
already linear for the tool: on the migration segment it asks $8$
goal-reachability queries at $n=1$ and $23$ at $n=6$, fewer than the projected
modular analysis' $48$. Nobody needs to decompose that. What is exponential is
computing an \emph{interface} --- enumerating the hazard-free terminations a
segment leaves behind --- and that is the like-for-like comparison, because it
is what a boundary contract costs. There the complete whole-system enumeration
asks $1536$ queries at $n=6$ against $48$ ($91.9$\,s against $0.222$\,s; we
report the counters, which reproduce exactly). The concrete interface asks
$2091$, more than the whole system, because it grows exponentially ($496$ points
at $n=6$) --- computing an unabstracted interface is worse than not abstracting
at all, and the projected one has at most three points at every $n$.

The win that survives is the practical one: re-checking after a change costs
whole-system analysis the full run ($23$ queries, $351$\,ms at $n=6$, growing
with $n$) and modular analysis the last segment alone ($8$ queries,
$36$\,ms, not growing with $n$).

\begin{table}[t]
\centering\small
\resizebox{\textwidth}{!}{%
\begin{tabular}{@{}rrrrrrr@{}}
\toprule
$n$ & whole, goal queries & exits & modular concrete, queries & interface & modular projected, queries & interface\\
\midrule
1 & 8 & 3 & 8 & 3 & 8 & 2\\
2 & 32 & 8 & 32 & 9 & 16 & 3\\
3 & 96 & 20 & 101 & 26 & 24 & 3\\
4 & 256 & 48 & 291 & 72 & 32 & 3\\
5 & 640 & 112 & 795 & 192 & 40 & 3\\
6 & 1536 & 256 & 2091 & 496 & 48 & 3\\
\bottomrule
\end{tabular}}
\caption{Chains of the migration segment, $k=2$. Re-check after replacing the
last segment: whole-system 28--351\,ms growing with $n$; modular 30--41\,ms,
not growing. The deployment segment gives the same shape at smaller scale
(2208\,ms vs 80\,ms at $n=6$).}
\label{tab:mod}
\end{table}

Composing through an interface rather than internals is a line of work in MPST
itself~\cite{gheri-oopsla23,barbanera-jlamp21}; what the budget adds is that the
interface carries a \emph{quantity} --- how much tolerance is left --- not only
a state. Theorem~\ref{thm:seq} does not make checking cheap; it makes it
\emph{decomposable}, and the interface form is what lets a coarse invariant
stand in for the concrete exits. Inferring good interfaces beyond the cone of
influence is the open engineering problem.

\subparagraph*{The base discipline.} The achievability checker the analysis is
built on was also evaluated on real skills. Those results concern the base
discipline, not severity, so we keep three facts and leave the data, the scripts
and the hand audits to the artifact. A token-free census of 157 skills from
public repositories refutes 14 even in their home runtime, and a hand audit finds
6 of those 14 genuine and the rest misextractions by the deterministic front end:
the refutations are sound for the pack the front end produced, not for the
document. In 136 agent runs over 16 documents, with every artifact verified by
recomputation, 68 of the 76 runs the checker certified produced a verified
artifact; but at realistic input size 5 of 16 refuted runs still computed a
correct artifact by hand, so a refusal is right about the skill as written, not
about the task.

\section{Repairs}\label{sec:repair}

When a branch is \Cat{} at the intended budget, four transformations make the
mistake affordable: \emph{guard} (insert a validation before the irreversible
branch), \emph{compensate} (add a recovery path), \emph{narrow} (remove the
branch, so the gate refuses it), and \emph{reorder} (move the irreversible
action after the validation). All have ancestors: amending asserted
choreographies~\cite{bocchi-lange-tuosto}, interactional
exceptions~\cite{carbone-concur08}, supervisory control and
shields~\cite{ramadge-wonham,bloem-tacas15,alshiekh-aaai18}, choreography
repair~\cite{basu-bultan,lanese-wwv13}. All four are mechanized (Coq:
\texttt{Repairs.v}). A \emph{validation} of $\psi$ is a capability that fires
exactly when $\psi$ holds and changes nothing --- a tool with precondition
$\psi$ and no effects.

\begin{theorem}[Narrow]\label{thm:narrow}
Removing branches from a choice node preserves the condition at the same
budget (\coqok{repair\_narrow\_sound}).
\end{theorem}

\begin{theorem}[Congruence, and the repairs off the root]\label{thm:congruence}
The condition is a congruence for one-hole contexts: if $G_2$ is $b$-safe
wherever $G_1$ is, then $C[G_2]$ is $b$-safe wherever $C[G_1]$ is, for a hole
under actions, markers and any branch of any choice
(\coqok{safeT\_congruence}). All four repairs therefore apply at nested
positions, which is where the tool applies them
(\coqok{repair\_narrow\_anywhere}, \coqok{repair\_guard\_anywhere},
\coqok{repair\_reorder\_anywhere}, \coqok{repair\_compensate\_anywhere}). The
price is that a repair whose soundness depends on the world --- guard needs an
inert abort world, compensate needs its recovery path safe at the exit points
--- must have that hypothesis at every world the context can reach, not only
at the root, and for the guard that is not free: a validation that validates
everywhere, into an idle abort world, forces the abort map to be idempotent on
its own image (\coqok{global\_abort\_is\_idempotent}). A context-wide guard
repair is therefore a repair with a single sink, not a fresh abort world per
failure. We state the constraint rather than hide it in a hypothesis; the
runtime we build (\texttt{Abort.v}) instantiates the root-level version.
\end{theorem}

How a validation is modelled decides whether the repair means anything. The
obvious model --- a capability that fires only when its predicate holds ---
makes the repaired protocol \emph{uninhabited} in exactly the worlds the repair
addresses: the conformance judgment's action rule asks for a successor, and
there is none, so the bridge would hold vacuously there. That is the shape of
defect that killed this discipline's predecessor (\S\ref{sec:whatchanged}).

A \emph{halted} world, one with no successors, only moves the vacuity one
constructor along: the guarded branch begins with an action, whose rule fails
there too, and no capability model here has one (\coqok{E2\_no\_halted}). An abort world is not a world where nothing is
enabled. It is a world where everything is enabled and nothing changes. We call
such a world \emph{idle}, and \emph{inert} if it is also hazard-free.

That condition is strong and we say plainly what it buys. \emph{Idle} is not
``the runtime answers and the answers are errors'': it is ``the session is
over''. $E_4$ (\texttt{Abort.v}) is the weaker reading --- the validation is a
validation at \emph{every} world, the abort flag records the failure, and every
other tool keeps working (\coqok{E4\_validates}, \coqok{E4\_no\_halted},
\coqok{E4\_not\_idle}) --- and there the same repair on the same protocol
fails: the guarded branch is not $1$-tolerant
(\coqok{Gguarded\_not\_1\_tolerant\_in\_E4}) and the reordered protocol is
tolerant at no budget (\coqok{reordered\_not\_tolerant\_in\_E4}). So the
repairs are about a gate that \emph{stops the session}, not one that reports a
failure; both discharge the misselected branch by that one step.

\begin{theorem}[A guarded branch is futile, not stuck]\label{thm:guardabort}
A validation of $\psi$ with abort map $\mathit{ab}$ is enabled in every world
(\coqok{validates\_ab\_enabled}). Nothing new is reachable from an idle world
(\coqok{reach\_idle}), so every protocol is safe from an inert one
(\coqok{safeT\_inert}); hence a misselected guarded branch, the goal failing at
the abort world, is \Fut{} --- the run is wasted, nothing is harmed --- rather
than stuck (\coqok{guard\_abort\_is\_futile}).
\end{theorem}

Every hypothesis of Theorem~\ref{thm:guardabort} is discharged concretely
(\texttt{Abort.v}). $E_2$ is a runtime with three capabilities --- verify,
purchase, and a validation of \emph{verified} --- all enabled in every world,
acting in a live world and behaving as the identity once an abort flag is set.
It has inert worlds (\coqok{E2\_inert\_abort}) and no halted ones, and its
validation satisfies the aborting model at every live world
(\coqok{E2\_validates}).

\begin{theorem}[Guard]\label{thm:guard}
Let $a$ validate $\psi$ with an inert abort world. A choice with the branch
$\ell[\psi].\,a.\,G_\ell$ added is $k$-safe iff the choice without it is and
$G_\ell$ is $k$-safe whenever $\psi$ holds (\coqok{repair\_guard\_exact\_ab}):
a misselected guarded branch diverts to the abort world, so it costs budget
but reaches nothing (\coqok{guard\_absorbs\_misselection\_ab}). The artifact
also carries the blocking-model versions (\coqok{repair\_guard\_exact}); they
are the same statement over a validation that cannot fire, and we do not rely
on them.
\end{theorem}

\begin{theorem}[Reorder]\label{thm:reorder}
If $\mathit{chk}$ never creates the hazard and commutes with $\mathit{irr}$
--- doing $\mathit{chk}$ then $\mathit{irr}$ re-serializes as $\mathit{irr}$ then
$\mathit{chk}$ with the same end world, as STRIPS effects on disjoint variables
do --- then $\mathit{chk}.\mathit{irr}.G$ is $k$-safe whenever
$\mathit{irr}.\mathit{chk}.G$ is (\coqok{repair\_reorder\_sound}). The two
protocols differ precisely at the point of no return: with $\psi$ false, an
inert abort world and one hazardous $\mathit{irr}$-successor, the original is
untypable and the reordered protocol is typable at every budget
(\coqok{repair\_reorder\_pnr\_ab}). Both are also characterized exactly in the
blocking model (\coqok{reorder\_original\_exact},
\coqok{reorder\_reordered\_exact}).
\end{theorem}

\begin{theorem}[Compensate]\label{thm:compensate}
Appending a recovery path $R$ to a residual $G_\ell$ is $k$-safe iff $R$ is
safe at every interface point of $G_\ell$ (\coqok{repair\_compensate\_sound},
an instance of Theorem~\ref{thm:seq}; the converse is not claimed), and it
turns a \Fut{} residual \Ben{} as soon as $R$ reaches the goal from one
interface point
(\coqok{repair\_compensate\_restores\_goal}).
\end{theorem}

Guard and reorder are exercised end to end on the booking protocol. Reorder: the
instance whose fast path purchases before verifying is not $1$-tolerant
(\coqok{Gbad2\_not\_1\_tolerant}); \emph{verify} and \emph{purchase} commute
(\coqok{commutes\_verify\_purchase}, proved without functional extensionality
by stating the capability context extensionally), and the reordered protocol
is tolerant at every $k$ (\coqok{Greordered\_is\_k\_tolerant}). Guard, over
$E_2$: the unguarded protocol is not $1$-tolerant
(\coqok{Gbad3\_not\_1\_tolerant}), prefixing the fast branch with the
validation makes it tolerant at every $k$
(\coqok{Gguarded\_is\_k\_tolerant}), a two-role session \emph{conforms} to the
repaired protocol (\coqok{Gguarded\_inhabited}) --- the obligation the
blocking model could not meet --- and so the bridge applies to it: every run
of that session is hazard-free at every budget
(\coqok{Gguarded\_bridge\_nonvacuous}).

\section{Related Work}\label{sec:related}

\subparagraph*{Runtime compliance monitoring.} The closest work by goal is
runtime compliance monitoring of language-model agents, of which
TRAC~\cite{trac-facct26} is the most developed: \textsc{ltl} constraints,
offline auditing, online monitoring, and predictive and intervening monitors
that sample continuations to preempt a violation. We share the question and differ on every
mechanism, in two ways. \emph{Prediction versus projection:} TRAC estimates a
coming violation by sampling a black box, where the guards fix the intended
branch and reachability the consequence of the others, so our prediction is
exact and free of any model of the agent, at the price of expressiveness.
\emph{The trusted seam:} TRAC's soundness bottoms out in learned labelers whose
temporal reasoning its own results show degrading with distance and constraint
count; ours in a deterministic checker. The two compose: TRAC's constraints can
supply $\Haz$, its monitor enforce what the condition certified. Choreographies
with first-class agent choice mapped to games~\cite{pact} own the quantitative
version of this framing; we are qualitative and static.

\subparagraph*{MPST with a fault model.} Crash-stop
failures~\cite{barwell-concur22,barwell-ecoop23,peters-forte22,viering-oopsla21,lebrun-esop23,hou-ecoop24},
affine and exceptional sessions~\cite{fowler-popl19,lagaillardie-ecoop22} and
protocol-induced recovery~\cite{neykova-cc17} model participants that stop;
misselection is one that continues, wrongly, and none classify the consequence
or track a goal. Asserted and refined global
types~\cite{bocchi-concur10,zhou-oopsla20,gheri-ecoop22,vassor-ecoop24} supply
our syntax; monitoring~\cite{bocchi-forte13} and blame~\cite{jia-popl16} treat
guard violation as an alarm rather than an object of analysis.

\subparagraph*{Budgeted deviation, elsewhere.} A budget on how many times an
assumption may be violated is not ours. Reactive synthesis has
$(k,b)$-resilience --- tolerate at most $k$ environment-assumption violations,
decided on a product with a decrementing counter~\cite{ehlers-hscc14} --- and
$k$-robustness for specifications~\cite{bloem-acta14}. Robustification computes
the largest deviation set a design survives~\cite{zhang-fse20,meiragoes-cav23},
which is our $\kstar$ in set currency, and planning has possibilistic
$K$-resilience~\cite{aineto-ecai23}; closest in fault model,
\emph{trembling-hand} \textsc{ltl}$_f$ planning~\cite{yu-ijcai24} assumes an
agent whose action selection is imprecise. We take the mechanism from this literature and
claim nothing about it. What is new is the carrier: the deviation is a
\emph{typed participant's own branch selection against an asserted guard}, the
budget is spent inside a session type, and Theorem~\ref{thm:bridge} carries it
through subject reduction to programs, where it distributes over roles
(Theorem~\ref{thm:distribute}). None of the works above types a program, and
none carries a quantity a participant spends across a reduction relation.

\subparagraph*{Bounded faults and dead ends.} Fault-tolerant
planning~\cite{domshlak-icaps13,aineto-ecai23} bounds action failures against
goal reachability; we bound selections against hazard reachability. Dead-end
detection~\cite{steinmetz-aaai16,hoffmann-ecai14},
undoability~\cite{daum-icaps16,morak-kr25} and excuse
generation~\cite{gobelbecker-icaps10} supply the point-of-no-return machinery;
reversibility in reinforcement-learning safety~\cite{turner-aies20} shares
the intuition without types.

\subparagraph*{Types and model checkers.} Naik and
Palsberg~\cite{naik-palsberg} and Kobayashi and Ong~\cite{kobayashi-ong} set
the template we instantiate. Bounded-model-checking of
fault-tolerant algorithms~\cite{stoilkovska-tacas19} is the verification-side
analogue of Theorem~\ref{thm:decide}.

\subparagraph*{Agents.} Formalisms for agents acting on irreversible external
state --- execution-fidelity invariants~\cite{yin-fidelity},
capability-containment proofs for skill manifests~\cite{metere-skills} ---
guarantee exactly-once effects or capability subsets; neither asks which wrong
choice is affordable. Runtime enforcement and provenance gating of tool
calls~\cite{agentspec,camel,fides,progent}, effect-typed agent
languages~\cite{etas}, network-enforced session types~\cite{nest-ecoop26} and
LTL trace monitoring~\cite{trac-facct26,agentltl} act on runs or in the host
language; we classify branches before any run exists.

\section{Limitations and Conclusion}\label{sec:conclusion}

\subparagraph*{A design that did not survive.}\label{sec:whatchanged} An earlier
version of this discipline typed per-role compliance --- trust modes, taint
propagation, staleness grades --- rather than severity. Its mechanized audit
found the typing rules vacuous on any protocol containing a capability, its
taint update blind to what a capability reads, and its loop condition satisfied
by cycles that never reset a grade; the audit is retained in the artifact. The
present design has no such premise --- deviation lives in the operational
semantics as a budget --- and that history is why this paper checks its own
theorems for vacuity (Theorems~\ref{thm:inhabited} and~\ref{thm:guardabort}).

\subparagraph*{Scope.} The development is mechanized for the head-move
semantics, which is the gate's; for ungated deployments
Theorem~\ref{thm:bystander} covers every permutation of independent nodes.
Asynchronous buffering is outside the synchronous discipline, and a reviewer
comparing scope against the mechanised asynchronous subject reduction
of~\cite{tirore-ecoop25} will be right to. Guardedness is a hypothesis of
progress alone, proved necessary rather than assumed
(Theorem~\ref{thm:guarded}). On the boolean fragment the tool is cross-checked
against a kernel extracted from the proof; the elaboration into its input
and the widened arithmetic fragment remain trusted; with environment choices,
recursion is mechanized for the hazard alone, and composition, the cone,
bystanders and repairs not at all. The
achievability verdict is about the skill \emph{as written}: an agent that
computes the artifact by hand is refuted with it, by design, and at realistic
sizes that route stays open for some tasks. The hazard is a declared or derived
input; \Ben{} inherits the checker's over-approximation, \Cat{} is a proof.
The benchmark is ours, its verdicts stated in advance but authored by us, and
only five of its seventeen protocols are in scope for the session-level
results; two of its ``repairs'' are hoists no theorem of ours licenses. The
real protocols are not ours but their world models are: we read the guards, the
preconditions and the irreversibility off policy text and tool code, the
expectations they were checked against are our hand reasoning, and the twelve
policy protocols rely on environment choices whose verdicts no kernel checks.
The live-agent experiment uses two models of one vendor, five runs per cell and
a simulated environment; with true facts it drew no error on our own protocols,
and all but one of its catastrophes come from a misinformed condition whose note we wrote. The guard and reorder repairs need a runtime whose failed
validation stops the session, not merely the call (\S\ref{sec:repair}).
Everything presumes the gated architecture.

\subparagraph*{Conclusion.} A participant that may never err may never act. The
question is not whether it will choose wrongly --- it will --- but whether the
protocol survives it: \emph{every affordable mistake is allowed, and every
unaffordable one is structurally unreachable.} On real policies live agents do
take the unaffordable ones, and the gate removed every one they took.

\bibliographystyle{plainurl}
\bibliography{main}
\end{document}
